\documentclass[11pt,letterpaper,openany]{book}

% ──────────────────────────────────────────────
% Packages
% ──────────────────────────────────────────────
\usepackage[margin=1.2in]{geometry}
\usepackage{amsmath, amssymb, amsthm}
\usepackage{mathtools}
\usepackage{bm}
\usepackage{tikz}
\usepackage{pgfplots}
\pgfplotsset{compat=1.18}
\usetikzlibrary{arrows.meta, decorations.markings, calc, patterns, positioning, 3d}
\usepackage{tcolorbox}
\tcbuselibrary{theorems, breakable, skins}
\usepackage{hyperref}
\usepackage{imakeidx}
\makeindex[intoc, title=Index]
\usepackage{enumitem}
\usepackage{fancyhdr}
\usepackage{titlesec}
\usepackage{epigraph}
\usepackage{xcolor}
\usepackage{booktabs}
\usepackage{listings}

% ──────────────────────────────────────────────
% Code Formatting (Listings)
% ──────────────────────────────────────────────
\definecolor{codegreen}{rgb}{0,0.6,0}
\definecolor{codegray}{rgb}{0.5,0.5,0.5}
\definecolor{codepurple}{rgb}{0.58,0,0.82}
\definecolor{backcolour}{rgb}{0.95,0.95,0.92}

\lstdefinestyle{mystyle}{
    backgroundcolor=\color{backcolour},
    commentstyle=\color{codegreen},
    keywordstyle=\color{magenta},
    numberstyle=\tiny\color{codegray},
    stringstyle=\color{codepurple},
    basicstyle=\ttfamily\footnotesize,
    breakatwhitespace=false,
    breaklines=true,
    captionpos=b,
    keepspaces=true,
    numbers=left,
    numbersep=5pt,
    showspaces=false,
    showstringspaces=false,
    showtabs=false,
    tabsize=2
}
\lstset{style=mystyle}

% ──────────────────────────────────────────────
% Colors
% ──────────────────────────────────────────────
\definecolor{chapblue}{RGB}{0, 70, 130}
\definecolor{accentred}{RGB}{180, 40, 40}
\definecolor{softgray}{RGB}{245, 245, 245}
\definecolor{trajgreen}{RGB}{30, 130, 60}
\definecolor{darkgold}{RGB}{160, 120, 20}
\definecolor{purplehaze}{RGB}{100, 40, 140}
\definecolor{biored}{RGB}{140, 20, 60}

% ──────────────────────────────────────────────
% Theorem environments
% ──────────────────────────────────────────────
\newtcbtheorem[number within=chapter]{principle}{Biomechanical Principle}{
  colback=chapblue!5, colframe=chapblue!80,
  fonttitle=\bfseries, separator sign={.~},
  breakable
}{princ}

\newtcbtheorem[number within=chapter]{clinicalbox}{Clinical Significance}{
  colback=biored!5, colframe=biored!80,
  fonttitle=\bfseries, separator sign={.~},
  breakable
}{clin}

\newtcbtheorem[number within=chapter]{example}{Example}{
  colback=softgray, colframe=black!40,
  fonttitle=\bfseries, separator sign={.~},
  breakable
}{ex}

\newtcbtheorem[number within=chapter]{definition}{Definition}{
  colback=darkgold!5, colframe=darkgold!80,
  fonttitle=\bfseries, separator sign={.~},
  breakable
}{def}

\newtcbtheorem[number within=chapter]{comparison}{Biology vs.\ Engineering}{
  colback=trajgreen!5, colframe=trajgreen!80,
  fonttitle=\bfseries, separator sign={.~},
  breakable
}{comp}







\newtcbtheorem[number within=chapter]{laymansbox}{In Plain Language}{
  colback=green!5!white, colframe=green!50!black,
  fonttitle=\bfseries, separator sign={.~},
  breakable
}{laymans}

\theoremstyle{plain}
\newtheorem{proposition}{Proposition}[chapter]
\newtheorem{theorem}{Theorem}[chapter]
\newtheorem{lemma}[theorem]{Lemma}

% ──────────────────────────────────────────────
% Chapter formatting
% ──────────────────────────────────────────────
\titleformat{\chapter}[display]
  {\normalfont\huge\bfseries\color{chapblue}}
  {\chaptertitlename\ \thechapter}{20pt}{\Huge}
\titlespacing*{\chapter}{0pt}{-20pt}{30pt}

% ──────────────────────────────────────────────
% Header/Footer
% ──────────────────────────────────────────────
\pagestyle{fancy}
\fancyhf{}
\fancyhead[LE,RO]{\thepage}
\fancyhead[RE]{\textit{Volume III: Biomechanics of Motion}}
\fancyhead[LO]{\textit{\nouppercase{\leftmark}}}
\renewcommand{\headrulewidth}{0.4pt}
\setlength{\headheight}{14.5pt}

% ──────────────────────────────────────────────
% Custom commands
% ──────────────────────────────────────────────
\newcommand{\state}{\bm{x}}
\newcommand{\control}{\bm{u}}
\newcommand{\traj}{\bm{\gamma}}
\newcommand{\manifold}{\mathcal{M}}
\newcommand{\configspace}{\mathcal{Q}}
\newcommand{\statespace}{\mathcal{X}}
\newcommand{\Reals}{\mathbb{R}}
\newcommand{\dd}{\mathrm{d}}
\newcommand{\dt}{\,\dd t}
\newcommand{\norm}[1]{\left\lVert #1 \right\rVert}
\newcommand{\inner}[2]{\left\langle #1,\, #2 \right\rangle}
\newcommand{\se}{\mathfrak{se}(3)}
\newcommand{\so}{\mathfrak{so}(3)}
\newcommand{\SE}{\mathrm{SE}(3)}
\newcommand{\SO}{\mathrm{SO}(3)}
\newcommand{\screw}{\mathcal{S}}
\newcommand{\twist}{\mathcal{V}}
\newcommand{\wrench}{\mathcal{F}}

% Biomechanics-specific commands
\newcommand{\muscle}{F_{\mathrm{m}}}
\newcommand{\tendon}{F_{\mathrm{t}}}
\newcommand{\activation}{a}
\newcommand{\momentarm}{r}
\newcommand{\pennation}{\alpha}

% ──────────────────────────────────────────────
% Title Page
% ──────────────────────────────────────────────
\title{%
  {\Huge\bfseries\color{chapblue} The Geometry of Motion}\\[0.4cm]
  {\Large Volume III: Biomechanics\\
  From Rigid Bodies to Biological Systems}\\[0.6cm]
  {\color{chapblue}\rule{\textwidth}{2pt}}
}
\date{}

\hypersetup{colorlinks=true, linkcolor=chapblue, filecolor=magenta, urlcolor=chapblue, citecolor=accentred}


\begin{document}

\maketitle
\thispagestyle{empty}
\cleardoublepage

\frontmatter
\chapter*{Preface}
\addcontentsline{toc}{chapter}{Preface}

The first two volumes of \emph{The Geometry of Motion} developed the mathematical language of
tangent spaces, contraction theory, and trajectory-centric control with the implicit assumption
that the systems under study are \emph{rigid bodies}---that joints are perfect revolutes, that
links are perfectly stiff, and that actuators produce torques on command. In biological systems,
\emph{none of these assumptions hold}.

This volume confronts reality. Biological tissues are viscoelastic, not rigid. Joints permit
coupled translations and rotations, not single-axis pivots. Muscles produce force through a
complex cascade from neural excitation through calcium dynamics to cross-bridge cycling, with
force depending on length, velocity, and activation history. Tendons store and release elastic
energy. Ligaments constrain motion in ways that change with loading rate.

And yet, the mathematical framework of Volumes~0--II is not abandoned---it is \emph{extended}.
The mass matrix $M(\bm{q})$ still exists, but it now includes the effective inertia of muscles
and the compliance of tendons. The configuration space $\mathcal{Q}$ is still a manifold, but
its geometry is richer because biological joints are not simple revolutes. The control input
$\bm{u}$ is no longer a torque but a neural excitation signal, and the mapping from excitation
to torque passes through the muscle model.

The philosophy remains unchanged: mathematics \emph{driven by physical reality}, with every
abstraction justified by measurement and every model validated against experiment. Where
engineering systems are designed to be simple, biological systems have evolved to be
\emph{robust}. Understanding why requires the tools of this volume.

\tableofcontents
\chapter*{Nomenclature}
\addcontentsline{toc}{chapter}{Nomenclature}

\section*{General Conventions}
\begin{itemize}
    \item Scalars are lowercase or Greek letters (e.g., $m, t, \theta$).
    \item Vectors are bold lowercase (e.g., $\bx, \bu, \bq$).
    \item Matrices are bold uppercase (e.g., $\bA, \bB, \bM, \bJ$).
    \item Expected values and sets are denoted by blackboard bold or script (e.g., $\mathbb{E}, \mathcal{S}$).
\end{itemize}

\section*{State and Kinematics}
\begin{description}
    \item[$\bq \in \Reals^n$] Joint configuration vector (generalized coordinates).
    \item[$\dot{\bq} \in \Reals^n$] Joint velocity vector (generalized velocities).
    \item[$\ddot{\bq} \in \Reals^n$] Joint acceleration vector.
    \item[$\bx \in \Reals^{nx}$] System state vector; commonly $\bx = [\bq^T, \dot{\bq}^T]^T$ for mechanical systems.
    \item[$\bu \in \Reals^m$] Control input vector (e.g., joint torques, muscle activations).
    \item[$\btau \in \Reals^n$] Generalized forces/torques acting on the joints.
    \item[$t \in \Reals$] Time.
\end{description}

\section*{Inertia and Dynamics}
\begin{description}
    \item[$\bM(\bq)$] Mass (inertia) matrix, symmetric positive definite.
    \item[$\bC(\bq, \dot{\bq})$] Coriolis and centrifugal force matrix.
    \item[$\bg(\bq)$] Gravity vector.
    \item[$\bJ(\bq)$] Jacobian matrix relating joint velocities to task-space velocities ($\dot{\bp} = \bJ(\bq)\dot{\bq}$).
\end{description}

\section*{Control and Optimization}
\begin{description}
    \item[$\bA_k = \frac{\partial f}{\partial \bx}$] Local Jacobian matrix of the state dynamics.
    \item[$\bB_k = \frac{\partial f}{\partial \bu}$] Local Jacobian matrix of the control input.
    \item[$V(\bx)$ or $\mathcal{V}$] Value function or Lyapunov function.
    \item[$\bQ$] State penalty matrix in LQR/DDP.
    \item[$\bR$] Control penalty matrix in LQR/DDP.
    \item[$\bK$] Feedback gain matrix ($\delta\bu = -\bK\delta\bx$).
    \item[$\traj$] A trajectory or flow, mapping time and initial conditions to states.
\end{description}

\section*{Biomechanics and Motor Control}
\begin{description}
    \item[$a_i \in [0, 1]$] Muscle activation level.
    \item[$\ell_i^M$] Muscle fiber length.
    \item[$v_i^M$] Muscle fiber contraction velocity.
    \item[$\ell_i^{MT}$] Musculotendon unit length.
    \item[$\bm{F}_{\text{muscle}}$] Force generated by muscles.
    \item[$\bW$] Muscle synergy matrix (weights).
    \item[$\bc(t)$] Muscle synergy activation vector over time.
\end{description}

\section*{Geometry and Manifolds}
\begin{description}
    \item[$SE(3)$] Special Euclidean group of rigid body motions.
    \item[$SO(3)$] Special Orthogonal group of 3D rotations.
    \item[$\mathfrak{se}(3)$] Lie algebra of $SE(3)$, space of spatial twists (velocities).
    \item[$\mathcal{M}$] A smooth manifold.
    \item[$T_{\bx}\mathcal{M}$] Tangent space at point $\bx$.
\end{description}

\cleardoublepage

\mainmatter

\chapter{How Biology Differs from Engineering}
\label{ch:bio-vs-eng}

\epigraph{%
  Nature does not build machines. It grows organisms.
  The engineer's model of the body is a useful fiction---but it is fiction.
}{--- Adapted from D'Arcy Wentworth Thompson}

\vspace{0.5cm}

\section{The Rigid-Body Assumption and Its Limits}

Every model in Volumes~0--II rested on a fundamental assumption: the moving
system consists of \emph{rigid bodies} connected by \emph{ideal joints}. A rigid body
has a fixed mass distribution; its shape does not change under load. An ideal revolute
joint constrains five of six relative degrees of freedom, permitting only rotation
about a single axis. These assumptions are extraordinarily productive for engineering
systems---robot arms, spacecraft, mechanisms---where components are designed to
approximate rigidity.

In biological systems, these assumptions fail at every level:

\begin{enumerate}[label=(\roman*)]
  \item \textbf{Bones are not rigid.} Long bones (femur, humerus, tibia) have measurable
        compliance under dynamic loading. A golf club shaft deflects by several centimeters
        and several degrees during a swing; a human femur deflects by millimeters during
        running. The deflection is small in absolute terms but can have significant effects
        on end-effector accuracy.

  \item \textbf{Joints are not revolutes.} The knee, for example, combines rolling,
        sliding, and rotation in a coupled manner that changes with flexion angle. The
        glenohumeral (shoulder) joint permits motion in three rotational axes but also
        allows several millimeters of translation. The intervertebral joints of the spine
        have six degrees of freedom each, all coupled through the disc and ligament
        mechanics.

  \item \textbf{Actuators are not torque sources.} Muscles produce force, not torque.
        Force is generated through the contraction of sarcomeres, transmitted through
        tendons, and converted to joint torque via moment arms that vary with joint angle.
        The force a muscle can produce depends on its length, its contraction velocity,
        and its activation history---dependencies that have no analog in electric motors.

  \item \textbf{Sensors are noisy, delayed, and multimodal.} Proprioceptors (muscle
        spindles, Golgi tendon organs, joint receptors) provide state information with
        significant noise and delays of 30--100~ms. The central nervous system must
        fuse information from proprioceptive, visual, and vestibular sources to estimate
        the body's state.

  \item \textbf{The system adapts.} Unlike a robot, the human body changes over
        time---muscles fatigue, connective tissue remodels, neural pathways strengthen
        or weaken. A complete model must account for adaptation on timescales from seconds
        (fatigue) to years (training).
\end{enumerate}

\begin{comparison}{The Fundamental Differences}{rigid-vs-bio}
\begin{center}
\renewcommand{\arraystretch}{1.3}
\begin{tabular}{@{}lll@{}}
\toprule
\textbf{Property} & \textbf{Engineering System} & \textbf{Biological System} \\
\midrule
Links & Rigid bodies & Viscoelastic, deformable \\
Joints & 1-DOF revolute/prismatic & Multi-DOF, coupled, with laxity \\
Actuators & Motors (torque on command) & Muscles (force via excitation-contraction) \\
Transmission & Gears, belts (rigid) & Tendons (elastic, energy-storing) \\
Sensors & Encoders (precise, fast) & Proprioceptors (noisy, delayed) \\
Parameters & Fixed (by design) & Variable (fatigue, adaptation, growth) \\
Redundancy & Minimal (cost-driven) & Massive (evolution-driven) \\
Failure mode & Brittle (sudden) & Graceful (degradation) \\
\bottomrule
\end{tabular}
\end{center}
\end{comparison}

These differences reflect fundamentally distinct design philosophies. Engineering systems are designed
for precise, repeatable control of pre-specified outputs. Biological systems are evolved for robustness,
adaptability, and efficiency under uncertain conditions. Neither approach is superior; rather, they
optimize for different objectives. Understanding which differences are essential for your particular
biomechanical question is crucial for successful modeling.

\begin{center}
\begin{tikzpicture}[scale=0.9, node distance=2.5cm]
  % Left side: Engineering system
  \node[rectangle, draw, thick, minimum width=1.8cm, minimum height=0.8cm, fill=blue!10] (eng-body) at (0, 0) {Rigid Body};
  \node[rectangle, draw, thick, minimum width=1.8cm, minimum height=0.8cm, fill=blue!10] (eng-joint) at (0, -1.5) {Ideal Joint};
  \node[rectangle, draw, thick, minimum width=1.8cm, minimum height=0.8cm, fill=blue!10] (eng-motor) at (0, -3) {Motor};

  \draw[->, thick] (eng-body) -- (eng-joint);
  \draw[->, thick] (eng-joint) -- (eng-motor);

  \node[above] at (0, 0.8) {\textbf{Engineering System}};
  \node[right] at (1.2, 0) {Fixed shape};
  \node[right] at (1.2, -1.5) {1-DOF};
  \node[right] at (1.2, -3) {Torque input};

  % Right side: Biological system
  \node[rectangle, draw, thick, minimum width=1.8cm, minimum height=0.8cm, fill=red!10] (bio-body) at (5, 0) {Bone + Tissue};
  \node[rectangle, draw, thick, minimum width=1.8cm, minimum height=0.8cm, fill=red!10] (bio-joint) at (5, -1.5) {Real Joint};
  \node[rectangle, draw, thick, minimum width=1.8cm, minimum height=0.8cm, fill=red!10] (bio-muscle) at (5, -3) {Muscle};

  \draw[->, thick] (bio-body) -- (bio-joint);
  \draw[->, thick] (bio-joint) -- (bio-muscle);

  \node[above] at (5, 0.8) {\textbf{Biological System}};
  \node[right] at (6.2, 0) {Deformable};
  \node[right] at (6.2, -1.5) {Multi-DOF, coupled};
  \node[right] at (6.2, -3) {Force + state};
\end{tikzpicture}
\end{center}

\section{Why the Differences Matter for Control}

The differences enumerated above are not merely academic. They fundamentally change
the control problem in ways that Volumes~I and~II must now accommodate.

\subsection{Compliance Creates Hidden Dynamics}

When a link is compliant, its deformation adds internal degrees of freedom to the system.
A golf club shaft, modeled as rigid, has one degree of freedom (the grip angle). Modeled
as a flexible beam, it has \emph{infinitely many} degrees of freedom, though practical
models truncate to the first 2--3 bending modes. The equation of motion becomes:
\begin{equation}
  \underbrace{M(\bm{q})\ddot{\bm{q}} + C(\bm{q},\dot{\bm{q}})\dot{\bm{q}} + g(\bm{q})}_{\text{rigid-body dynamics}}
  + \underbrace{K_{\text{stiff}} \bm{q}_{\text{flex}} + D_{\text{damp}} \dot{\bm{q}}_{\text{flex}}}_{\text{flexibility terms}}
  = \bm{\tau},
  \label{eq:ch1:flexible-eom}
\end{equation}
where $\bm{q}_{\text{flex}}$ represents the modal coordinates of the flexible modes,
$K_{\text{stiff}}$ is the stiffness matrix, and $D_{\text{damp}}$ is the structural
damping matrix.

\begin{principle}{The Timescale Separation Principle}{timescale}
  In many biological systems, the flexible dynamics are \emph{fast} relative to the
  rigid-body dynamics. If the natural frequency of the flexible modes is much higher
  than the bandwidth of the controller, the flexibility can be modeled as a
  quasi-static compliance rather than a dynamic mode. Specifically, if the $k$-th
  flexible mode has natural frequency $\omega_k$ and the controller bandwidth is
  $\omega_c$, then:
  \begin{itemize}
    \item $\omega_k \gg \omega_c$: quasi-static (flexibility ``follows'' the rigid motion)
    \item $\omega_k \approx \omega_c$: resonance risk (flexibility must be modeled dynamically)
    \item $\omega_k \ll \omega_c$: the system is effectively another rigid body (the
          flexible mode is so slow it acts like a separate DOF)
  \end{itemize}
\end{principle}

The timescale separation principle is critical for model selection. For different body
regions and different movement speeds, the applicability of the rigid-body approximation varies.
For the human body during locomotion, empirical measurements suggest that long bones have natural
frequencies sufficiently high that the rigid-body assumption is justified. For the spine and
soft tissues, particularly during rapid movements, the separation is less clean, and dynamic
flexibility effects may contribute meaningfully to the system behavior. The choice of whether to
include flexibility as explicit state variables or treat it as a quasi-static effect should be
guided by both the timescale analysis above and by quantitative comparison of model predictions
to experimental data.

\subsection{Muscle Force Production Is State-Dependent}

The most profound difference between engineered and biological actuators is that muscle
force depends on the \emph{state of the muscle itself}. An electric motor produces torque
proportional to current, regardless of shaft angle or angular velocity (to first
approximation). A muscle's maximum force depends on:

\begin{enumerate}
  \item \textbf{Muscle length} (force-length relationship): maximum isometric force
        occurs at an intermediate ``optimal'' length $\ell_0$, with force decreasing
        at shorter and longer lengths.
  \item \textbf{Contraction velocity} (force-velocity relationship): concentric
        (shortening) contractions produce less force at higher velocities; eccentric
        (lengthening) contractions can produce forces exceeding the maximum isometric
        force by 50--80\%.
  \item \textbf{Activation level}: the fraction of motor units recruited and their
        firing rates, which follows its own first-order dynamics with time constants
        of 10--50~ms for activation and 40--200~ms for deactivation.
\end{enumerate}

The combined force production model (Hill model, Chapter~3) is:
\begin{equation}
  F_{\text{muscle}} = a(t) \cdot f_L(\tilde{\ell}) \cdot f_V(\tilde{v}) \cdot F_0
  + f_{\text{PE}}(\tilde{\ell}),
  \label{eq:ch1:hill-overview}
\end{equation}
where $a(t) \in [0,1]$ is the activation level, $f_L$ and $f_V$ are the normalized
force-length and force-velocity curves, $F_0$ is the maximum isometric force,
$\tilde{\ell} = \ell/\ell_0$ and $\tilde{v} = v/v_0$ are normalized length and velocity,
and $f_{\text{PE}}$ is the passive elastic force.

\begin{principle}{Force-Velocity Asymmetry in Multi-Joint Tasks}{fv-asymmetry-task}
  In any task where joints undergo different phases of contraction (some concentric,
  others eccentric), the force-velocity relationship creates an asymmetry in joint
  torque capacity. This asymmetry can be quantified through the Hill force-velocity
  function applied to each muscle group. The engineering consequence is that a model
  assuming constant or task-invariant joint torque limits systematically misrepresents
  the instantaneous torque capacity of biological joints, particularly during rapid
  multi-joint movements.
\end{principle}

\subsection{Redundancy Is a Feature, Not a Bug}

The human body has approximately 630 skeletal muscles controlling approximately 240
joint degrees of freedom. Even for a simple task like reaching for a cup (which
constrains only 6 degrees of freedom---3 position + 3 orientation of the hand),
there are vastly more muscles and joints than constraints. This \emph{motor redundancy}
means that infinitely many muscle activation patterns can produce the same hand
trajectory.

In engineering, redundancy is minimized to reduce cost. In biology, redundancy is
\emph{the solution to robustness}:

\begin{itemize}
  \item If one muscle is fatigued, others can compensate.
  \item If a joint is injured, the task can be accomplished through alternative
        kinematic chains.
  \item Variability in execution---far from being noise to be eliminated---allows
        the system to explore the space of solutions and find robust attractors
        (see Volume~IV for the motor control implications).
\end{itemize}

This connects directly to the \emph{uncontrolled manifold hypothesis}
\cite{ScholzSchoner1999}: task-irrelevant variability is not suppressed by the
nervous system, only task-relevant variability is. The biological control
system does not solve the full redundancy problem; it projects the problem
onto the task-relevant subspace.

\section{Modeling Strategies}

Given the complexity of biological systems, modelers must choose their level of
abstraction carefully. Three broad strategies exist, listed in order of increasing
fidelity:

\subsection{Strategy 1: Rigid-Body with Torque Actuators}

The simplest approach: treat biological segments as rigid bodies and joints as
ideal revolutes, with ``joint torques'' as the control input. This is the model
assumed throughout Volumes~I--II.

\textbf{Advantages:}
\begin{itemize}
  \item Mathematically tractable (standard Lagrangian mechanics)
  \item Sufficient for gross motion analysis
  \item Compatible with all control-theoretic tools from Volumes~I--II
\end{itemize}

\textbf{Limitations:}
\begin{itemize}
  \item Cannot predict muscle activation patterns
  \item Cannot account for muscle force-length-velocity constraints
  \item Cannot predict injury mechanisms
  \item Cannot capture bi-articular muscle effects
\end{itemize}

\textbf{When to use:} Trajectory optimization, gross motion planning, initial
controller design.

\subsection{Strategy 2: Rigid-Body with Muscle Actuators}

Replace torque actuators with Hill-type muscle models (Chapter~3). Joint torques
become \emph{outputs} of the model, computed from muscle forces and moment arms:
\begin{equation}
  \tau_j = \sum_{i \in \mathcal{M}_j} r_{ij}(\bm{q}) \cdot F_{\text{muscle},i},
  \label{eq:ch1:muscle-torque}
\end{equation}
where $\mathcal{M}_j$ is the set of muscles crossing joint $j$ and $r_{ij}(\bm{q})$
is the moment arm of muscle $i$ about joint $j$.

\textbf{Advantages:}
\begin{itemize}
  \item Predicts muscle activation patterns
  \item Captures force-length-velocity constraints
  \item Accounts for bi-articular coupling
  \item Can predict co-contraction and joint stiffness modulation
\end{itemize}

\textbf{Limitations:}
\begin{itemize}
  \item Significant increase in model complexity (typically 2--3 states per muscle)
  \item Muscle parameters are difficult to measure in vivo
  \item Still assumes rigid segments and ideal joints
\end{itemize}

\textbf{When to use:} Detailed motion analysis, muscle force prediction, clinical
biomechanics, prosthetic design.

\subsection{Strategy 3: Flexible Multi-Body with Muscle Actuators}

The most complete approach: flexible bodies (beam or FEM models) with muscle
actuators and realistic joint models. This is the domain of software platforms
like OpenSim and AnyBody.

\textbf{Advantages:}
\begin{itemize}
  \item Highest fidelity
  \item Can predict stress distributions in bones and soft tissues
  \item Can capture coupled joint kinematics
\end{itemize}

\textbf{Limitations:}
\begin{itemize}
  \item Computationally expensive
  \item Many parameters, often poorly known
  \item Too complex for real-time control applications
\end{itemize}

\textbf{When to use:} Injury prediction, surgical planning, detailed structural
analysis.

\begin{lstlisting}[language=Python, caption={Choosing a modeling strategy based on application requirements.}]
import numpy as np

def select_modeling_strategy(
    need_muscle_forces: bool,
    need_tissue_stress: bool,
    real_time_required: bool,
    n_dof: int
) -> str:
    """Select the appropriate biomechanical modeling strategy.

    Parameters
    ----------
    need_muscle_forces : bool
        Whether individual muscle forces are needed.
    need_tissue_stress : bool
        Whether stress/strain in tissues is needed.
    real_time_required : bool
        Whether the model must run in real-time.
    n_dof : int
        Number of degrees of freedom in the model.

    Returns
    -------
    str
        Recommended modeling strategy.
    """
    if need_tissue_stress:
        return "Strategy 3: Flexible multi-body + muscles (OpenSim/FEM)"
    elif need_muscle_forces:
        if real_time_required and n_dof > 20:
            return ("Strategy 2 (reduced): Muscle model with "
                    "synergy-based dimensionality reduction")
        return "Strategy 2: Rigid-body + Hill-type muscles (OpenSim)"
    else:
        if real_time_required:
            return "Strategy 1: Rigid-body + torque actuators (fastest)"
        return ("Strategy 1 or 2 depending on available data "
                "and analysis goals")
\end{lstlisting}

\section{A Quantitative Tour of the Human Body}

Before diving into the detailed modeling chapters, we present a quantitative overview
of the key biomechanical properties that any model must capture. All values in this
section are drawn from published anthropometric studies and are cited in the relevant
subsections.

\subsection{Segment Inertial Properties}

The human body can be decomposed into 15--20 segments for modeling purposes. The
standard segmentation and mass fractions are based on cadaver studies, most notably
those of Dempster (1955), updated by de~Leva (1996) \cite{deLeva1996}.

\begin{center}
\renewcommand{\arraystretch}{1.2}
\begin{tabular}{@{}lccc@{}}
\toprule
\textbf{Segment} & \textbf{Mass (\% body)} & \textbf{Length (cm, 180cm male)} & \textbf{CoM (\% proximal)} \\
\midrule
Head + neck & 8.1 & 25.4 & 50.0 \\
Trunk & 49.7 & 53.0 & 44.9 \\
Upper arm & 2.7 & 28.2 & 57.7 \\
Forearm & 1.6 & 26.9 & 45.7 \\
Hand & 0.6 & 19.0 & 36.2 \\
Thigh & 14.2 & 42.2 & 40.9 \\
Shank & 4.3 & 43.4 & 43.4 \\
Foot & 1.4 & 25.8 & 44.2 \\
\bottomrule
\end{tabular}
\end{center}

\begin{lstlisting}[language=Python, caption={Computing segment inertial properties from body measurements.}]
import numpy as np
from dataclasses import dataclass

@dataclass
class SegmentProperties:
    """Inertial properties of a body segment."""
    name: str
    mass: float          # kg
    length: float        # m
    com_proximal: float  # fraction from proximal end
    I_com: float         # moment of inertia about CoM, kg*m^2

def de_leva_male_segments(
    body_mass: float,
    body_height: float
) -> list[SegmentProperties]:
    """Compute segment properties using de Leva (1996) regression.

    Parameters
    ----------
    body_mass : float
        Total body mass in kg.
    body_height : float
        Total body height in m.

    Returns
    -------
    list[SegmentProperties]
        List of segment inertial properties.

    References
    ----------
    de Leva, P. (1996). Adjustments to Zatsiorsky-Seluyanov's
    segment inertia parameters. J. Biomechanics, 29(9), 1223-1230.
    """
    # Mass fractions, length fractions, CoM fractions, RoG fractions
    # (proximal reference, male data from de Leva 1996 Table 4)
    data = {
        'head_neck':  (0.0810, 0.1414, 0.5002, 0.3030),
        'trunk':      (0.4970, 0.2946, 0.4486, 0.3271),
        'upper_arm':  (0.0271, 0.1566, 0.5772, 0.2850),
        'forearm':    (0.0162, 0.1494, 0.4574, 0.2760),
        'hand':       (0.0061, 0.1057, 0.3624, 0.2880),
        'thigh':      (0.1416, 0.2321, 0.4095, 0.3290),
        'shank':      (0.0433, 0.2411, 0.4340, 0.2550),
        'foot':       (0.0137, 0.1527, 0.4415, 0.2570),
    }

    segments = []
    for name, (mass_frac, len_frac, com_frac, rog_frac) in data.items():
        seg_mass = mass_frac * body_mass
        seg_length = len_frac * body_height
        seg_com = com_frac
        # Moment of inertia: I = m * (rog * length)^2
        seg_I = seg_mass * (rog_frac * seg_length) ** 2
        segments.append(SegmentProperties(
            name=name, mass=seg_mass, length=seg_length,
            com_proximal=seg_com, I_com=seg_I
        ))

    return segments


# Example usage
if __name__ == "__main__":
    segments = de_leva_male_segments(body_mass=80.0, body_height=1.80)
    print(f"{'Segment':<15} {'Mass (kg)':>10} {'Length (m)':>10} "
          f"{'I_com (kg*m^2)':>14}")
    print("-" * 52)
    for seg in segments:
        print(f"{seg.name:<15} {seg.mass:>10.2f} {seg.length:>10.3f} "
              f"{seg.I_com:>14.4f}")
\end{lstlisting}

\subsection{Joint Ranges of Motion}

Each joint has characteristic ranges of motion that define the configuration space
$\mathcal{Q}$ of the body. Unlike robot joints, which can be modeled as
$q_i \in [\underline{q}_i, \overline{q}_i]$, biological joints have \emph{soft limits}:
resistance increases gradually as the joint approaches its anatomical limit, governed
by the elasticity of ligaments and joint capsule.

\subsection{Muscle Properties}

The human body contains approximately 630 skeletal muscles. Each is characterized by
several anatomical and physiological parameters that determine its force production capacity.
The standard parameters are:

\begin{itemize}
  \item \textbf{Maximum isometric force} $F_0$: the force a muscle produces at optimal
        length with full activation and zero velocity. Values vary widely by muscle size,
        composition, and fiber type.
  \item \textbf{Optimal fiber length} $\ell_0$: the length at which maximum isometric force
        is produced, determined by the sarcomere length ($\sim$2.6 $\mu$m per sarcomere at
        $\ell_0$).
  \item \textbf{Maximum contraction velocity} $v_0$: the velocity at which a muscle produces
        zero force under concentric contraction, measured in fiber lengths per second (typically
        $0.5$--$3$ fiber lengths per second for slow-twitch fibers, $5$--$15$ for fast-twitch).
  \item \textbf{Tendon slack length} $\ell_T^s$: the length of the tendon at which it begins to
        transmit force (below this length, the tendon is slack and carries no load).
  \item \textbf{Pennation angle} $\alpha_0$: the angle between muscle fiber direction and the
        tendon line of action at optimal fiber length. Ranges from approximately $0^\circ$ for
        parallel-fibered muscles to $30^\circ$ or more for highly pennate muscles.
\end{itemize}

These properties have been systematically measured and catalogued through:
\begin{enumerate}
  \item \cite{Zajac1989}: foundational review of muscle mechanics and parameter estimation
  \item \cite{Delp2007}: comprehensive OpenSim musculoskeletal model with measured parameters for major muscles
  \item \cite{Arnold2010}: detailed muscle morphometry from imaging and cadaver studies
\end{enumerate}

Values for specific muscles are provided in software platforms (OpenSim 4.x, AnyBody) and should
be obtained from published measurements rather than estimated.

\section{Connecting to the Control Framework}

Despite the added complexity, the control framework of Volumes~I and~II remains
applicable---with modifications. The key insight is:

\begin{principle}{The Extended Affine Structure}{affine-bio}
  A musculoskeletal system can be written in control-affine form as
  \begin{equation}
    \dot{\state} = f(\state) + G(\state) \bm{u},
    \label{eq:ch1:affine-bio}
  \end{equation}
  where the state $\state$ now includes muscle states (activation, fiber length)
  in addition to joint coordinates and velocities, the drift field $f$ includes
  passive muscle and tendon forces, and $G(\state)$ maps neural excitation signals
  $\bm{u}$ to state derivatives through the muscle activation and contraction dynamics.
  The dimensions expand---from $2n$ states with $n$ joints and $m$ torque actuators
  to $2n + 2p$ states with $p$ muscles---but the \emph{structure} of the problem
  (drift + control) is preserved.
\end{principle}

\begin{remark}{Dimensional Expansion of the State Space}{expanded_manifold}
  Recall from Volumes~0 and~I that the state manifold $\mathcal{M}$ has dimension $2n$ for a system
  with $n$ joint degrees of freedom (tracking positions and velocities). When muscles are included,
  the state space expands: for each of $p$ muscles, we add two dimensions to track activation $a_i$
  and fiber length $\ell_i^M$, yielding a state space of dimension $2n + 2p$. The system still evolves
  according to $\dot{\state} = f(\state) + G(\state) \bm{u}$, but the ambient manifold is now much
  higher-dimensional. All differential geometry tools (Lie brackets, contraction metrics, etc.)\ from
  Volumes~I--II apply without modification, though computational complexity scales with dimension.
\end{remark}

This means:
\begin{itemize}
  \item The tangent-space analysis of Volume~I applies to the expanded state space
  \item Contraction theory (Volume~I, Chapter~4) can assess stability of the
        muscle-driven system
  \item Trajectory optimization (Volume~I, Chapter~5) extends to find optimal
        neural excitation patterns
  \item The counterfactual analysis (Volume~I, Chapter~7) can distinguish which
        aspects of muscle force production contribute to the motion
\end{itemize}

The challenge is \emph{dimensionality}: a 20-DOF body with 100 muscles has a
state space of dimension $2 \times 20 + 2 \times 100 = 240$. This is precisely
the ``curse of dimensionality'' that Volume~IV addresses through the lens of
motor control and neural architecture.

\section*{Exercises}

\begin{enumerate}
  \item \textbf{Segment Properties.}
        Using the \texttt{de\_leva\_male\_segments()} function, compute the total
        mass and center of mass of a 75-kg, 1.75-m male in the anatomical position.
        Verify that the total mass sums to the body mass.

  \item \textbf{Compliance Effects.}
        A golf club shaft has a first bending mode at approximately 5~Hz. If the
        downswing takes 0.25~s, how many oscillation cycles of the shaft occur during
        the downswing? Does the timescale separation principle (Principle~\ref{princ:timescale})
        justify neglecting shaft flexibility? Discuss.

  \item \textbf{Redundancy Quantification.}
        A planar model of the arm has 3 joints (shoulder, elbow, wrist) and 6 muscles.
        If the task is to place the fingertip at a given $(x, y)$ position, what is the
        dimension of the null space? How many independent muscle patterns produce the same
        fingertip position?

  \item \textbf{Force-Velocity Impact.}
        During a golf downswing, the wrist extensors lengthen at approximately 5 fiber
        lengths per second. Using the Hill force-velocity relationship (Equation~\ref{eq:ch1:hill-overview}),
        estimate the ratio of eccentric force to maximum isometric force at this velocity.
        Compare this with the force available during a slow concentric contraction.

  \item \textbf{Modeling Strategy Selection.}
        For each of the following applications, select the appropriate modeling strategy
        and justify your choice:
        \begin{enumerate}[label=(\alph*)]
          \item Real-time control of a prosthetic knee during walking
          \item Predicting ACL stress during a cutting maneuver in soccer
          \item Optimizing a sprinter's start to minimize 10-m time
          \item Understanding how the golf swing generates clubhead speed
        \end{enumerate}

  \item \textbf{(Programming.)} Implement a function that computes the total gravitational
        potential energy of a multi-segment body given segment properties and joint angles.
        Test it for a 2-segment (upper arm + forearm) model in different configurations.

  \item \textbf{(Programming.)} Extend the segment properties code to compute the full
        $6 \times 6$ spatial inertia matrix (in the Park and Lynch spatial algebra
        notation) for each segment. Verify that the matrices are positive definite.

  \item \textbf{(Research.)} Compare the de~Leva (1996) segment parameters with at least
        one other source (e.g., Winter 2009, Dempster 1955). Quantify the differences in
        mass fractions and discuss the implications for inverse dynamics calculations.
\end{enumerate}

\chapter{Musculoskeletal Modeling Conventions}
\label{ch:msk-conventions}

\epigraph{%
  In biomechanics, the first source of error is not the mathematics
  but the coordinate system.
}{--- ISB Standardization Committee}

\section{Anatomical Reference Frames}

Biomechanical analysis requires a standardized language for describing body position and
motion. The International Society of Biomechanics (ISB) has published recommendations
for joint coordinate systems \cite{Wu2002, Wu2005} that serve as the international
standard.

\begin{center}
\begin{tikzpicture}[scale=1.0, >=Stealth]
  % Global frame
  \draw[thick, ->, red] (0, 0, 0) -- (2, 0, 0) node[right] {$\hat{\bm{x}}_G$ (anterior)};
  \draw[thick, ->, green] (0, 0, 0) -- (0, 2, 0) node[above] {$\hat{\bm{y}}_G$ (superior)};
  \draw[thick, ->, blue] (0, 0, 0) -- (0, 0, 2) node[below right] {$\hat{\bm{z}}_G$ (lateral)};

  \node at (-0.3, -0.3, 0) {$\{G\}$};

  % Segment frame (offset)
  \draw[thick, ->, red!60] (3, 1, 0) -- (4.5, 1.2, 0) node[right, font=\small] {$\hat{\bm{x}}_S$};
  \draw[thick, ->, green!60] (3, 1, 0) -- (3, 2.5, 0) node[above, font=\small] {$\hat{\bm{y}}_S$};
  \draw[thick, ->, blue!60] (3, 1, 0) -- (3, 1, 1.5) node[right, font=\small] {$\hat{\bm{z}}_S$};

  \node at (2.7, 0.5, 0) {$\{S\}$};

  % Bony landmarks on segment
  \filldraw[fill=black] (3, 1, 0) circle (3pt) node[below, font=\tiny] {Origin};
  \filldraw[fill=gray] (3.5, 2.8, 0) circle (2pt);
  \filldraw[fill=gray] (3.8, 2.5, 0) circle (2pt);

  \node[font=\tiny, gray] at (4.2, 2.9) {Landmarks};
\end{tikzpicture}
\end{center}

\subsection{The Anatomical Position}

All references to body position begin from the \textbf{anatomical position}: the subject
stands upright, feet together, arms at sides, palms facing forward. This defines the
\emph{zero configuration} from which all joint angles are measured.

\begin{definition}{Global Reference Frame}{global-frame}
  The global (laboratory) reference frame $\{G\}$ is defined as:
  \begin{itemize}
    \item $\hat{\bm{x}}_G$: anterior (forward)
    \item $\hat{\bm{y}}_G$: superior (upward)
    \item $\hat{\bm{z}}_G$: lateral right
  \end{itemize}
  This is a right-handed coordinate system. Note that different communities use
  different conventions: biomechanics typically uses $Y$-up, while robotics often
  uses $Z$-up. \textbf{Always verify the convention before combining data from
  different sources.}
\end{definition}

\subsection{Segment Reference Frames}

Each body segment has a local reference frame attached to it, defined by bony landmarks
that can be palpated or identified from medical imaging. The ISB recommendations specify
these landmarks for each segment \cite{Wu2002, Wu2005}.

\begin{example}{The Thigh Reference Frame}{thigh-frame}
  The ISB-recommended thigh reference frame $\{T\}$ is defined by:
  \begin{itemize}
    \item \textbf{Origin}: midpoint between medial and lateral femoral epicondyles
    \item $\hat{\bm{y}}_T$: from origin to hip joint center (proximal direction)
    \item $\hat{\bm{z}}_T$: perpendicular to the plane formed by the two epicondyles
          and the hip center, pointing laterally
    \item $\hat{\bm{x}}_T$: cross product $\hat{\bm{y}}_T \times \hat{\bm{z}}_T$
          (anterior direction)
  \end{itemize}
  The hip joint center is estimated from the pelvis landmarks using regression
  equations (e.g., \cite{Harrington2007}).
\end{example}

\subsection{Connection to Screw Theory}

In the language of Volume~0, each segment frame $\{S_i\}$ is related to the global
frame $\{G\}$ by a homogeneous transformation matrix $T_{GS_i} \in \SE$:
\begin{equation}
  T_{GS_i} = \begin{bmatrix} R_{GS_i} & \bm{p}_{GS_i} \\ \bm{0}^T & 1 \end{bmatrix},
  \label{eq:ch2:homogeneous}
\end{equation}
where $R_{GS_i} \in \SO$ is the rotation matrix and $\bm{p}_{GS_i} \in \Reals^3$ is
the position of the segment origin in the global frame.

The \emph{relative configuration} of segment $j$ with respect to proximal segment $i$
is:
\begin{equation}
  T_{S_i S_j} = T_{GS_i}^{-1} T_{GS_j}.
  \label{eq:ch2:relative-config}
\end{equation}
Joint angles are extracted from this relative transformation.

\section{Joint Coordinate Systems}

The ISB recommends decomposing the relative rotation $R_{S_i S_j}$ into three
angles using a specific Euler angle sequence for each joint. The choice of
sequence is \emph{not arbitrary}---it is selected to align with the clinical
definitions of flexion-extension, abduction-adduction, and internal-external
rotation.

\subsection{The Grood--Suntay Convention for the Knee}

The knee joint coordinate system, originally proposed by Grood and Suntay (1983)
\cite{GroodSuntay1983}, decomposes knee motion into:

\begin{enumerate}
  \item \textbf{Flexion-Extension} ($\alpha$): rotation about the femoral medio-lateral
        axis (the ``fixed'' axis of the femur)
  \item \textbf{Abduction-Adduction} ($\beta$): rotation about the ``floating'' axis
        (perpendicular to both the femoral and tibial body-fixed axes)
  \item \textbf{Internal-External Rotation} ($\gamma$): rotation about the tibial
        longitudinal axis (the ``body-fixed'' axis of the tibia)
\end{enumerate}

This is a ZXY Euler angle sequence when expressed in the Grood--Suntay convention.
The key insight is that the floating axis is neither fixed in the femur nor in the
tibia---it is defined by the cross product of the other two axes.

\begin{lstlisting}[language=Python, caption={Computing knee joint angles using the Grood-Suntay convention.}]
import numpy as np
from typing import Tuple

def grood_suntay_knee_angles(
    R_femur: np.ndarray,
    R_tibia: np.ndarray
) -> Tuple[float, float, float]:
    """Compute knee joint angles using Grood-Suntay convention.

    Parameters
    ----------
    R_femur : np.ndarray, shape (3, 3)
        Rotation matrix of femur segment frame in global coordinates.
    R_tibia : np.ndarray, shape (3, 3)
        Rotation matrix of tibia segment frame in global coordinates.

    Returns
    -------
    Tuple[float, float, float]
        (flexion_deg, adduction_deg, internal_rotation_deg)

    References
    ----------
    Grood, E.S. and Suntay, W.J. (1983). A joint coordinate system
    for the clinical description of three-dimensional motions.
    J. Biomechanical Engineering, 105(2), 136-144.
    """
    # Femoral medio-lateral axis (e1) - flexion axis
    e1 = R_femur[:, 2]  # lateral axis of femur

    # Tibial longitudinal axis (e3) - rotation axis
    e3 = R_tibia[:, 1]  # proximal axis of tibia

    # Floating axis (e2) - perpendicular to both
    e2 = np.cross(e3, e1)
    e2 = e2 / np.linalg.norm(e2)

    # Flexion-extension angle
    # Angle between femur anterior axis and floating axis
    femur_anterior = R_femur[:, 0]
    flexion = np.arctan2(
        np.dot(e2, R_femur[:, 1]),
        np.dot(e2, femur_anterior)
    )

    # Abduction-adduction angle
    # Angle between e1 and e3 minus 90 degrees
    sin_abd = -np.dot(e1, e3)
    cos_abd = np.linalg.norm(np.cross(e1, e3))
    adduction = np.arcsin(np.clip(sin_abd, -1.0, 1.0))

    # Internal-external rotation
    # Angle between floating axis and tibia anterior axis
    tibia_anterior = R_tibia[:, 0]
    internal_rot = np.arctan2(
        np.dot(e2, R_tibia[:, 2]),
        np.dot(e2, tibia_anterior)
    )

    return (
        np.degrees(flexion),
        np.degrees(adduction),
        np.degrees(internal_rot)
    )
\end{lstlisting}

\subsection{Multi-Representation Joint Angles}

Following the multi-representation philosophy of this series, we express joint
rotations in all four standard representations:

\textbf{1. Euler Angles (ISB Convention):}
\begin{equation}
  R_{ij} = R_z(\alpha) \, R_x(\beta) \, R_y(\gamma),
  \label{eq:ch2:euler-isb}
\end{equation}
where $\alpha, \beta, \gamma$ are the clinical angles (flexion, adduction, rotation).

\textbf{2. Rotation Matrix (SO(3)):}
\begin{equation}
  R_{ij} = \exp([\bm{\omega}]_\times \theta) \in \SO,
  \label{eq:ch2:rotmat}
\end{equation}
which avoids gimbal lock but requires 9 parameters with 6 constraints.

\textbf{3. Unit Quaternion:}
\begin{equation}
  \bm{q} = (\cos\tfrac{\theta}{2},\; \bm{\omega}\sin\tfrac{\theta}{2}),
  \quad \norm{\bm{q}} = 1,
  \label{eq:ch2:quaternion}
\end{equation}
which is singularity-free and computationally efficient.

\textbf{4. Screw Axis (Primary):}
\begin{equation}
  T_{ij} = e^{[\screw]\theta} \in \SE,
  \quad \screw = (\bm{\omega}, \bm{v}) \in \Reals^6,
  \label{eq:ch2:screw}
\end{equation}
which captures the coupled rotation-translation of biological joints naturally.

\begin{comparison}{Why Screw Axis Theory Is Natural for Biological Joints}{screw-bio}
  Unlike engineering revolute joints (pure rotation about a fixed axis), biological
  joints exhibit \emph{coupled} rotation and translation. The knee slides anteriorly
  as it flexes; the glenohumeral joint translates inferiorly during abduction.
  Screw axis theory captures this coupling naturally: every rigid-body motion is a
  rotation about and translation along a screw axis. The instantaneous screw axis
  (ISA) of a biological joint moves through space as the joint moves, unlike the
  fixed axis of an ideal revolute.

  This is precisely why Park and Lynch's \emph{Modern Robotics} framework is
  preferred: it handles coupled rotation-translation without special cases.
\end{comparison}

\section{Body Segment Parameters}

\subsection{Regression Equations}

The most widely used body segment parameter (BSP) data comes from:
\begin{enumerate}
  \item \textbf{Dempster (1955)}: cadaver study of 8 elderly males. Still widely cited
        but limited demographic range.
  \item \textbf{Zatsiorsky \& Seluyanov (1983, 1990)}: gamma-ray scanning of 100 living
        young males. More representative but uses different segment definitions.
  \item \textbf{de Leva (1996)}: adjustment of Zatsiorsky data to standard segment
        definitions. The current gold standard for adult male and female BSP.
  \item \textbf{Pavol et al.\ (2002)}: elderly adult data.
  \item \textbf{Jensen (1986, 1989)}: pediatric data.
\end{enumerate}

\subsection{Computing the Mass Matrix}

Given segment inertial properties, the mass matrix $M(\bm{q}) \in \Reals^{n \times n}$
of the multi-body system is computed using the composite rigid-body algorithm
(Volume~0, Chapter~10). In the body-fixed frame of segment $i$, the spatial
inertia is:
\begin{equation}
  \mathcal{G}_i = \begin{bmatrix} I_i & m_i [\bm{c}_i]_\times^T \\
  m_i [\bm{c}_i]_\times & m_i \mathbb{I}_3 \end{bmatrix} \in \Reals^{6 \times 6},
  \label{eq:ch2:spatial-inertia}
\end{equation}
where $I_i$ is the $3 \times 3$ rotational inertia about the segment frame origin,
$m_i$ is the segment mass, $\bm{c}_i$ is the center of mass position in the segment
frame, and $[\cdot]_\times$ denotes the skew-symmetric matrix.

\section*{Exercises}

\begin{enumerate}
  \item \textbf{Coordinate System Verification.}
        Verify that the Grood--Suntay knee joint coordinate system produces zero
        angles when the femur and tibia reference frames are aligned. Test with
        known rotations.

  \item \textbf{Euler Angle Singularity.}
        For the ZXY Euler angle sequence used in the knee, identify the gimbal lock
        configuration. Is this configuration ever reached during normal gait?
        During deep squatting?

  \item \textbf{Multi-Representation Conversion.}
        Write Python functions to convert between all four representations
        (Euler ZXY, rotation matrix, quaternion, screw axis) for a knee joint angle
        of $60^\circ$ flexion, $5^\circ$ adduction, $10^\circ$ internal rotation.

  \item \textbf{BSP Sensitivity.}
        Using the de~Leva segment properties, compute the knee extension torque
        required to hold the lower leg horizontal (isometric). Repeat using the
        Dempster data. What is the percentage difference? Discuss the implications for
        inverse dynamics accuracy.

  \item \textbf{(Programming.)} Implement the spatial inertia matrix computation
        for a 3-segment leg model (thigh, shank, foot). Verify positive-definiteness.

  \item \textbf{(Programming.)} Using motion capture data (or simulated data),
        compute knee joint angles from segment orientations using the Grood--Suntay
        code provided. Plot flexion angle over one gait cycle.
\end{enumerate}

\chapter{Muscle Models}
\label{ch:muscle-models}

\epigraph{%
  A muscle is not a motor. It is a spring with a memory,
  a damper with an agenda, and a force generator with a personality.
}{}

\section{The Hill-Type Muscle Model}

The Hill-type muscle model, originating from A.V.\ Hill's Nobel Prize-winning work on
muscle thermodynamics \cite{Hill1938}, is the standard lumped-parameter model used in
biomechanical simulation. It represents the musculotendon unit as a combination of
contractile, elastic, and damping elements.

\subsection{Model Architecture}

The standard Hill-type musculotendon model consists of three elements:

\begin{enumerate}
  \item \textbf{Contractile Element (CE)}: represents the active force-generating
        capacity of the muscle fibers. Its force depends on activation $a(t)$, fiber
        length $\ell^M$, and fiber velocity $v^M = \dot{\ell}^M$.
  \item \textbf{Parallel Elastic Element (PE)}: represents the passive elastic
        properties of the muscle belly (titin filaments, connective tissue). Produces
        force only when the muscle is stretched beyond its resting length.
  \item \textbf{Series Elastic Element (SE)}: represents the tendon, which is in
        series with the contractile element. Modeled as a nonlinear spring with a
        slack length $\ell_T^s$.
\end{enumerate}

\begin{center}
\begin{tikzpicture}[scale=1.0, >=Stealth]
  % Tendon (series elastic)
  \draw[thick] (0, 0) -- (1, 0);
  \draw[thick, decorate, decoration={zigzag, segment length=6, amplitude=3}]
    (1, 0) -- (3, 0) node[midway, above=8pt] {\small SE (tendon)};
  \draw[thick] (3, 0) -- (4, 0);

  % Junction point
  \filldraw (4, 0) circle (2pt);

  % Contractile element (top branch)
  \draw[thick] (4, 0) -- (4, 0.8) -- (5, 0.8);
  \draw[thick, fill=blue!10] (5, 0.4) rectangle (7, 1.2);
  \node at (6, 0.8) {\small CE};
  \draw[thick] (7, 0.8) -- (8, 0.8) -- (8, 0);

  % Parallel elastic element (bottom branch)
  \draw[thick] (4, 0) -- (4, -0.8) -- (5, -0.8);
  \draw[thick, decorate, decoration={zigzag, segment length=6, amplitude=3}]
    (5, -0.8) -- (7, -0.8) node[midway, below=8pt] {\small PE};
  \draw[thick] (7, -0.8) -- (8, -0.8) -- (8, 0);

  % Attachment points
  \draw[thick] (8, 0) -- (9, 0);
  \filldraw (0, 0) circle (2pt);
  \filldraw (9, 0) circle (2pt);

  % Labels
  \node[below] at (0, -0.1) {\small Origin};
  \node[below] at (9, -0.1) {\small Insertion};

  % Pennation angle
  \draw[thick, dashed, gray] (5, 0) -- (5, 0.4);
  \draw[gray, ->] (5, 0) arc (270:300:0.4);
  \node[gray, right] at (5.3, -0.2) {\small $\alpha$};
\end{tikzpicture}
\end{center}

\subsection{The Force-Length Relationship}

The isometric (zero velocity) force of the contractile element depends on the
normalized fiber length $\tilde{\ell} = \ell^M / \ell_0^M$:

\begin{definition}{Active Force-Length Curve}{fl-curve}
  The normalized active force-length relationship is:
  \begin{equation}
    f_L(\tilde{\ell}) = \exp\!\left(-\frac{(\tilde{\ell} - 1)^2}{\gamma}\right),
    \label{eq:ch3:force-length}
  \end{equation}
  where $\gamma \approx 0.45$ controls the width of the curve. Maximum force ($f_L = 1$)
  occurs at optimal fiber length ($\tilde{\ell} = 1$). Force decreases for both shorter
  and longer fiber lengths, reaching near-zero at approximately $\tilde{\ell} = 0.5$
  and $\tilde{\ell} = 1.5$.
\end{definition}

The physical basis for the force-length relationship is the sliding filament theory
\cite{HuxleySimmons1971}: at the optimal length, the maximum number of actin-myosin
cross-bridges can form. At shorter lengths, filament overlap reduces the number of
available binding sites. At longer lengths, the filaments pull apart.

\begin{lstlisting}[language=Python, caption={Hill-type muscle model implementation.}]
import numpy as np
from dataclasses import dataclass
from typing import Tuple

@dataclass
class MuscleParameters:
    """Parameters for a Hill-type musculotendon model."""
    F0: float          # Maximum isometric force (N)
    l0_M: float        # Optimal fiber length (m)
    v0: float          # Maximum contraction velocity (m/s)
    l_T_slack: float   # Tendon slack length (m)
    alpha0: float      # Pennation angle at optimal length (rad)
    gamma: float = 0.45  # Force-length curve width parameter
    Af: float = 0.25   # Force-velocity shape factor
    kPE: float = 4.0   # Passive elastic stiffness parameter
    e0: float = 0.6    # Passive elastic strain at F0

    @property
    def v_max(self) -> float:
        """Maximum shortening velocity in fiber lengths/s."""
        return self.v0 / self.l0_M


def force_length_active(l_tilde: float, gamma: float = 0.45) -> float:
    """Active force-length relationship.

    Parameters
    ----------
    l_tilde : float
        Normalized fiber length (l_M / l0_M).
    gamma : float
        Width parameter of the Gaussian curve.

    Returns
    -------
    float
        Normalized active force (0 to 1).

    References
    ----------
    Thelen, D.G. (2003). Adjustment of muscle mechanics model
    parameters to simulate dynamic contractions in older adults.
    J. Biomechanical Engineering, 125(1), 70-77.
    """
    return np.exp(-((l_tilde - 1.0) ** 2) / gamma)


def force_length_passive(l_tilde: float, kPE: float = 4.0,
                         e0: float = 0.6) -> float:
    """Passive force-length relationship (exponential model).

    Parameters
    ----------
    l_tilde : float
        Normalized fiber length.
    kPE : float
        Stiffness parameter.
    e0 : float
        Strain at which passive force equals F0.

    Returns
    -------
    float
        Normalized passive force.
    """
    if l_tilde <= 1.0:
        return 0.0
    return (np.exp(kPE * (l_tilde - 1.0) / e0 - kPE) - 1.0) / \
           (np.exp(kPE) - 1.0)


def force_velocity(v_tilde: float, a: float, Af: float = 0.25,
                   Flen: float = 1.4) -> float:
    """Force-velocity relationship (Thelen 2003 model).

    Parameters
    ----------
    v_tilde : float
        Normalized fiber velocity (v_M / v_max).
        Negative = shortening (concentric).
    a : float
        Activation level (0 to 1).
    Af : float
        Shape factor for concentric portion.
    Flen : float
        Maximum normalized eccentric force.

    Returns
    -------
    float
        Normalized force multiplier.
    """
    if v_tilde <= 0:
        # Concentric (shortening)
        return (1.0 + v_tilde) / (1.0 - v_tilde / (Af * a + 0.01))
    else:
        # Eccentric (lengthening)
        return (1.0 + v_tilde * Flen / (0.04 + a))  / \
               (1.0 + v_tilde / (0.04 + a))


def muscle_force(
    activation: float,
    l_M: float,
    v_M: float,
    params: MuscleParameters
) -> Tuple[float, float, float]:
    """Compute total musculotendon force.

    Parameters
    ----------
    activation : float
        Neural activation level (0 to 1).
    l_M : float
        Current muscle fiber length (m).
    v_M : float
        Current muscle fiber velocity (m/s, negative = shortening).
    params : MuscleParameters
        Muscle parameters.

    Returns
    -------
    Tuple[float, float, float]
        (total_force, active_force, passive_force) in Newtons.
    """
    # Normalize
    l_tilde = l_M / params.l0_M
    v_tilde = v_M / params.v0

    # Force components
    f_active = (activation
                * force_length_active(l_tilde, params.gamma)
                * force_velocity(v_tilde, activation, params.Af))
    f_passive = force_length_passive(l_tilde, params.kPE, params.e0)

    # Pennation angle (changes with fiber length)
    cos_alpha = np.cos(params.alpha0) if l_tilde > 0.1 else 1.0

    # Total force along tendon
    F_active = f_active * params.F0 * cos_alpha
    F_passive = f_passive * params.F0 * cos_alpha
    F_total = F_active + F_passive

    return F_total, F_active, F_passive


# Example: Biceps brachii
biceps = MuscleParameters(
    F0=624.3,        # N (Arnold et al. 2010)
    l0_M=0.1157,     # m
    v0=10 * 0.1157,  # 10 fiber lengths/s
    l_T_slack=0.2723, # m
    alpha0=0.0        # parallel-fibered
)

# Compute force at optimal length, half activation, isometric
F_total, F_active, F_passive = muscle_force(
    activation=0.5,
    l_M=biceps.l0_M,  # optimal length
    v_M=0.0,           # isometric
    params=biceps
)
print(f"Biceps at 50% activation, optimal length, isometric:")
print(f"  Active force:  {F_active:.1f} N")
print(f"  Passive force: {F_passive:.1f} N")
print(f"  Total force:   {F_total:.1f} N")
\end{lstlisting}

\subsection{The Force-Velocity Relationship}

The second critical property of the contractile element is the force-velocity
relationship, first characterized by Hill (1938) \cite{Hill1938}.

\begin{principle}{The Force-Velocity Asymmetry}{fv-asymmetry}
  \begin{itemize}
    \item \textbf{Concentric (shortening) contraction}: force \emph{decreases} with
          increasing shortening velocity. At maximum shortening velocity $v_0$, the
          muscle produces zero force. The relationship is approximately hyperbolic:
          $(F + a)(v + b) = (F_0 + a) b$.
    \item \textbf{Eccentric (lengthening) contraction}: force \emph{increases} beyond
          the maximum isometric force, typically reaching $1.4$--$1.8 \times F_0$.
          This is why muscles are ``stronger'' when resisting a load than when lifting it.
    \item \textbf{Isometric ($v = 0$)}: force equals the product of activation and
          the force-length relationship.
  \end{itemize}
  The asymmetry between concentric and eccentric force production is unique to
  biological actuators. No common engineering actuator exhibits this behavior.
\end{principle}

\subsection{Activation Dynamics}

The neural command from the central nervous system does not instantly produce
muscle force. The activation dynamics models the excitation-contraction coupling:
\begin{equation}
  \dot{a}(t) = \frac{u(t) - a(t)}{\tau(u, a)},
  \label{eq:ch3:activation-dynamics}
\end{equation}
where $u(t) \in [0, 1]$ is the neural excitation, $a(t) \in [0, 1]$ is the
activation level, and $\tau$ is a time constant that differs for activation
and deactivation:
\begin{equation}
  \tau(u, a) = \begin{cases}
    \tau_{\text{act}}(0.5 + 1.5a) & \text{if } u > a \text{ (activation)}, \\
    \tau_{\text{deact}} / (0.5 + 1.5a) & \text{if } u \leq a \text{ (deactivation)}.
  \end{cases}
  \label{eq:ch3:tau-activation}
\end{equation}
Typical values: $\tau_{\text{act}} = 10$--$15$~ms, $\tau_{\text{deact}} = 40$--$60$~ms
\cite{Winters1995}. The asymmetry between activation and deactivation time constants
has important consequences for motor control: it is much easier to ``turn on'' a
muscle quickly than to ``turn it off.''

\section{Tendon Mechanics}

The tendon connects the muscle belly to bone. It is modeled as a nonlinear elastic
element with a resting (slack) length $\ell_T^s$:
\begin{equation}
  F_T(\varepsilon_T) = \begin{cases}
    0 & \varepsilon_T \leq 0, \\
    F_0 \cdot \frac{e^{k_T \varepsilon_T} - 1}{e^{k_T \varepsilon_{0T}} - 1} &
    \varepsilon_T > 0,
  \end{cases}
  \label{eq:ch3:tendon-force}
\end{equation}
where $\varepsilon_T = (\ell_T - \ell_T^s) / \ell_T^s$ is the tendon strain and
$\varepsilon_{0T} \approx 0.033$ is the strain at maximum isometric force. The
tendon stiffness parameter $k_T$ is typically chosen so that the tendon strain is
approximately 3.3\% at $F_0$ \cite{Zajac1989}.

\begin{clinicalbox}{Energy Storage in Tendons}{tendon-energy}
  Tendons are not merely passive connectors---they are \emph{energy storage devices}.
  During running, the Achilles tendon stores approximately 35~J of elastic energy
  per stride, which is returned during push-off. This reduces the metabolic cost
  of running by approximately 50\% compared to a hypothetical system with rigid
  tendons \cite{Alexander1991}.

  For the golf swing, the tendons of the wrist and forearm muscles store energy
  during the early downswing (when the wrist is cocked) and release it during
  the uncocking phase, contributing to clubhead speed.
\end{clinicalbox}

\section{Musculotendon Equilibrium}

The muscle fiber and tendon must satisfy force equilibrium at all times. Assuming
the muscle fiber acts at a pennation angle $\alpha$ to the tendon line of action:
\begin{equation}
  F_T = (F_{\text{CE}} + F_{\text{PE}}) \cos\alpha,
  \label{eq:ch3:mt-equilibrium}
\end{equation}
where:
\begin{align}
  F_{\text{CE}} &= a \cdot f_L(\tilde{\ell}^M) \cdot f_V(\tilde{v}^M) \cdot F_0, \\
  F_{\text{PE}} &= f_{\text{PE}}(\tilde{\ell}^M) \cdot F_0.
\end{align}

The musculotendon length $\ell^{MT}$ is the sum of tendon length and the
projection of fiber length:
\begin{equation}
  \ell^{MT} = \ell_T + \ell^M \cos\alpha.
  \label{eq:ch3:mt-length}
\end{equation}

Given a known musculotendon length (from the skeletal geometry) and activation,
the equilibrium condition \eqref{eq:ch3:mt-equilibrium} and the constraint
\eqref{eq:ch3:mt-length} determine the fiber length and velocity.

\section{The Complete Hill Model as a Dynamical System}

Combining all elements, the Hill muscle model can be expressed as a dynamical
system with two state variables per muscle:
\begin{align}
  \dot{a} &= \frac{u - a}{\tau(u, a)}, \label{eq:ch3:state-activation} \\
  \dot{\ell}^M &= v^M(a, \ell^M, \ell^{MT}), \label{eq:ch3:state-fiber}
\end{align}
where $v^M$ is obtained by inverting the force-velocity relationship at the
equilibrium fiber force. The input is the neural excitation $u(t) \in [0, 1]$
and the ``disturbance'' is the musculotendon path length $\ell^{MT}(\bm{q})$,
which depends on the joint configuration.

\begin{principle}{Muscle as a Nonlinear Control System}{muscle-control}
  Each muscle is itself a nonlinear dynamical system with:
  \begin{itemize}
    \item \textbf{State}: $(a, \ell^M) \in [0,1] \times \Reals_{>0}$
    \item \textbf{Input}: neural excitation $u \in [0,1]$
    \item \textbf{Output}: tendon force $F_T$
    \item \textbf{Disturbance}: musculotendon path length $\ell^{MT}(\bm{q})$
  \end{itemize}
  The muscle transforms a scalar neural command into a scalar force, but the
  transformation is highly nonlinear, state-dependent, and history-dependent.
  The entire musculoskeletal system is a cascade: neural commands $\to$ muscle
  dynamics $\to$ tendon forces $\to$ joint torques $\to$ skeletal dynamics.
\end{principle}

\section*{Exercises}

\begin{enumerate}
  \item \textbf{Force-Length Curve.}
        Plot the active force-length curve for $\gamma \in \{0.3, 0.45, 0.6\}$.
        For each, determine the ``useful range'' (the interval of fiber lengths
        where $f_L > 0.5$). How does the width parameter affect the operating range?

  \item \textbf{Force-Velocity Curve.}
        Plot the complete force-velocity curve (both concentric and eccentric) for
        the biceps parameters given in the code listing. Mark the maximum eccentric
        force and the maximum shortening velocity.

  \item \textbf{Activation Dynamics.}
        Simulate the activation dynamics for a square-wave neural excitation input
        ($u = 1$ for 200~ms, then $u = 0$). Plot the activation response with
        $\tau_{\text{act}} = 15$~ms and $\tau_{\text{deact}} = 50$~ms. How long does
        it take for activation to reach 95\% of maximum? How long for deactivation?

  \item \textbf{Tendon Compliance.}
        For the Achilles tendon (slack length $\ell_T^s = 0.25$~m, max force 4000~N),
        compute the tendon stretch at 50\%, 75\%, and 100\% of maximum force. Does
        the tendon behave approximately linearly over this range?

  \item \textbf{Isometric Force Surface.}
        Create a 3D surface plot of muscle force as a function of normalized fiber
        length ($\tilde{\ell}$) and activation ($a$) at zero velocity (isometric).
        Identify the point of maximum force production.

  \item \textbf{(Programming.)} Implement a complete musculotendon simulation: given
        a prescribed musculotendon length trajectory $\ell^{MT}(t)$ and neural
        excitation $u(t)$, integrate the state equations \eqref{eq:ch3:state-activation}
        and \eqref{eq:ch3:state-fiber} and plot the resulting tendon force $F_T(t)$.

  \item \textbf{(Programming.)} Implement a two-muscle antagonist pair (e.g., biceps
        and triceps) acting on a single-DOF elbow joint. Given co-contraction
        ($u_{\text{biceps}} = u_{\text{triceps}} = 0.5$), compute the net joint torque
        and the joint stiffness. Verify that co-contraction increases stiffness without
        changing the equilibrium position.

  \item \textbf{(Research.)} Compare the Thelen (2003) muscle model with the
        Millard et al.\ (2013) model used in OpenSim 4.x. What are the key differences
        in the force-velocity and passive force models? Which is more physiologically
        accurate?
\end{enumerate}

\chapter{Joint Kinematics and Soft Tissue}
\label{ch:joint-kinematics}

\epigraph{%
  The knee is not a hinge. The shoulder is not a ball-and-socket.
  These are useful fictions that biology tolerates but does not enforce.
}{}

\section{Beyond Ideal Joints}

Engineering joints are designed to constrain motion to specific degrees of freedom:
a revolute joint permits one rotation, a prismatic joint permits one translation.
Biological joints achieve constraint through a fundamentally different mechanism:
the interaction of articular surfaces, ligaments, joint capsule, and muscle forces.
This section develops the kinematics of biological joints using the screw theory
framework of Volume~0, extended to handle the additional complexity.

\subsection{The Instantaneous Screw Axis}

Every rigid-body motion can be decomposed into a rotation about and translation along
a \textbf{screw axis} (Chasles' theorem, Volume~0, Chapter~5). For an engineering
revolute joint, this screw axis is fixed in both bodies. For a biological joint,
the instantaneous screw axis (ISA) \emph{moves} as the joint configuration changes.

\begin{definition}{Instantaneous Screw Axis of a Biological Joint}{isa-bio}
  Given the relative twist $\twist_{ij} = (\bm{\omega}, \bm{v}) \in \Reals^6$ between
  adjacent segments $i$ and $j$, the instantaneous screw axis has:
  \begin{itemize}
    \item Direction: $\hat{\bm{s}} = \bm{\omega} / \norm{\bm{\omega}}$
    \item Pitch: $h = \bm{\omega} \cdot \bm{v} / \norm{\bm{\omega}}^2$
    \item Location: $\bm{q} = \bm{\omega} \times \bm{v} / \norm{\bm{\omega}}^2$
  \end{itemize}
  For a pure rotation ($h = 0$), the ISA passes through $\bm{q}$ in direction $\hat{\bm{s}}$.
  Tracking how $\bm{q}$ and $\hat{\bm{s}}$ change with joint angle reveals the coupled
  kinematics of the biological joint.
\end{definition}

\subsection{The Knee: Rolling, Sliding, and Coupled Rotation}

The tibiofemoral joint exemplifies the complexity of biological kinematics. During
flexion from $0^\circ$ to $90^\circ$:
\begin{enumerate}
  \item The femoral condyles \textbf{roll} on the tibial plateaus (like a wheel on a road)
  \item They simultaneously \textbf{slide} posteriorly (the contact point moves backward)
  \item The tibia undergoes coupled \textbf{internal rotation} of approximately
        $15^\circ$ (the ``screw-home'' mechanism) \cite{walker1972quantitative,komistek2003invivo}
  \item The contact area shifts, changing the effective moment arm
\end{enumerate}

This coupled motion means the knee ISA migrates superiorly and posteriorly during
flexion. The four-bar linkage model \cite{OConnor1989} provides a simplified but
insightful mechanical analog:

\begin{lstlisting}[language=Python, caption={Instantaneous screw axis computation for a biological joint.}]
import numpy as np
from typing import Tuple

def compute_isa(
    omega: np.ndarray,
    v: np.ndarray
) -> Tuple[np.ndarray, np.ndarray, float]:
    """Compute the instantaneous screw axis from a twist.

    Parameters
    ----------
    omega : np.ndarray, shape (3,)
        Angular velocity vector.
    v : np.ndarray, shape (3,)
        Linear velocity of the body-fixed frame origin.

    Returns
    -------
    Tuple[np.ndarray, np.ndarray, float]
        (point_on_axis, axis_direction, pitch)
    """
    omega_norm = np.linalg.norm(omega)
    if omega_norm < 1e-10:
        # Pure translation
        direction = v / np.linalg.norm(v)
        return np.zeros(3), direction, np.inf

    direction = omega / omega_norm
    pitch = np.dot(omega, v) / omega_norm**2
    point = np.cross(omega, v) / omega_norm**2

    return point, direction, pitch


def knee_isa_trajectory(
    flexion_angles: np.ndarray,
    roll_slide_ratio: float = 0.6
) -> np.ndarray:
    """Estimate knee ISA position during flexion.

    Simple model: ISA migrates posteriorly and superiorly
    with flexion, governed by the roll-to-slide ratio.

    Parameters
    ----------
    flexion_angles : np.ndarray
        Array of flexion angles in degrees.
    roll_slide_ratio : float
        Ratio of rolling to total contact displacement.
        Pure rolling = 1.0, pure sliding = 0.0. A value of
        approximately 0.6 is consistent with the in-vitro
        tibiofemoral measurements of Walker et al. (1972)
        [walker1972quantitative].

    Returns
    -------
    np.ndarray, shape (N, 3)
        ISA positions in the tibial reference frame.
    """
    # Approximate femoral condyle radius (m); representative
    # adult value from Walker et al. (1972) [walker1972quantitative]
    # and standard knee-anatomy references [nigg2007biomechanics].
    R_condyle = 0.030
    isa_positions = np.zeros((len(flexion_angles), 3))

    for i, theta in enumerate(np.radians(flexion_angles)):
        # Posterior migration due to rolling
        posterior = R_condyle * theta * roll_slide_ratio
        # Superior migration (ISA moves up with flexion)
        superior = R_condyle * (1.0 - np.cos(theta * 0.5))

        isa_positions[i] = [-posterior, superior, 0.0]

    return isa_positions
\end{lstlisting}

\subsection{The Shoulder: The Most Complex Joint}

The glenohumeral joint has the largest range of motion of any joint in the body,
achieved at the expense of stability. The ``ball-and-socket'' description is a
gross simplification: the humeral head translates $2$--$4$~mm during abduction
\cite{harryman1990translation}, and the instantaneous center of rotation shifts
continuously.

The shoulder complex involves four joints:
\begin{enumerate}
  \item \textbf{Glenohumeral} (GH): primary ball-and-socket, 3 DOF rotation + $\sim$3~mm translation \cite{harryman1990translation}
  \item \textbf{Acromioclavicular} (AC): 3 DOF rotation, small range
  \item \textbf{Sternoclavicular} (SC): 3 DOF rotation, small range
  \item \textbf{Scapulothoracic} (ST): not a true synovial joint; the scapula glides on
        the thorax, constrained as a \emph{surface contact} rather than a pin joint
\end{enumerate}

The scapulothoracic ``joint'' is particularly interesting from a mathematical
perspective: it is a \emph{holonomic constraint} that the scapula must remain on
the thoracic surface, effectively acting as a 2-DOF constraint (the scapula can
translate on the surface but cannot lift off).

\section{Ligament Models}

Ligaments constrain joint motion by resisting excessive displacement. They are
modeled as nonlinear springs with:
\begin{equation}
  F_{\text{lig}}(\varepsilon) = \begin{cases}
    0 & \varepsilon \leq 0 \text{ (slack)}, \\
    k_1 \varepsilon^2 / (4\varepsilon_1) & 0 < \varepsilon \leq 2\varepsilon_1 \text{ (toe region)}, \\
    k_2(\varepsilon - \varepsilon_1) & \varepsilon > 2\varepsilon_1 \text{ (linear region)},
  \end{cases}
  \label{eq:ch4:ligament-force}
\end{equation}
where $\varepsilon = (\ell - \ell_0) / \ell_0$ is the strain, $\ell_0$ is the
reference length, $\varepsilon_1 \approx 0.03$ is the transition strain, and
$k_1, k_2$ are stiffness parameters \cite{Blankevoort1991}.

\section{Spinal Mechanics}

The spine presents a unique modeling challenge: 24 vertebrae connected by intervertebral
discs, each joint having 6 coupled degrees of freedom. The functional spinal unit (FSU)
consisting of two adjacent vertebrae and their connecting disc is the basic mechanical
unit.

The intervertebral disc consists of:
\begin{itemize}
  \item \textbf{Nucleus pulposus}: a gel-like center that distributes compressive loads
  \item \textbf{Annulus fibrosus}: concentric rings of collagen fibers that resist
        torsion and bending
\end{itemize}

A common simplified model treats each FSU as a 6-DOF spring:
\begin{equation}
  \bm{F}_{\text{disc}} = K_{\text{disc}} \cdot \bm{\delta},
  \label{eq:ch4:disc-stiffness}
\end{equation}
where $K_{\text{disc}} \in \Reals^{6 \times 6}$ is the stiffness matrix and
$\bm{\delta} \in \Reals^6$ is the generalized displacement (3 translations + 3 rotations).

\section*{Exercises}

\begin{enumerate}
  \item \textbf{ISA Computation.}
        For a knee flexing at $60^\circ/\text{s}$ with $2$~mm/s of posterior translation,
        compute the ISA direction, pitch, and location.

  \item \textbf{Screw-Home Mechanism.}
        Plot the coupled internal rotation of the tibia as a function of knee flexion
        using the model: $\gamma = 15^\circ \cdot (1 - \cos(\theta / 2))$ for
        $\theta \in [0, 90^\circ]$.

  \item \textbf{Shoulder Kinematics.}
        Using screw axis theory, express the configuration of the hand (end-effector)
        as a product of exponentials through the SC, AC, GH, and elbow joints.

  \item \textbf{(Programming.)} Implement the ligament force model and plot the
        force-strain curve. Identify the toe region and the linear region.

  \item \textbf{(Programming.)} Create a 3D visualization of the knee ISA trajectory
        during flexion from $0^\circ$ to $120^\circ$.
\end{enumerate}

\chapter{Multi-Body Dynamics of Biological Chains}
\label{ch:multibody-bio}

\epigraph{%
  The human body is the most complex open kinematic chain in nature.
  Analyzing it with the tools of rigid-body dynamics is both
  necessary and insufficient.
}{}

\section{Open-Chain vs.\ Closed-Chain Biomechanics}

Most human movement involves open kinematic chains (the limbs) connected to a floating
base (the pelvis/trunk). However, many functional activities create temporary closed
chains: the stance phase of gait (foot on ground), pushing against a wall, or the golf
swing when both hands grip the club.

\subsection{Equations of Motion for Biological Chains}

The equations of motion for an $n$-DOF musculoskeletal system take the standard
Lagrangian form:
\begin{equation}
  M(\bm{q})\ddot{\bm{q}} + C(\bm{q}, \dot{\bm{q}})\dot{\bm{q}} + g(\bm{q})
  = R(\bm{q}) \bm{F}_{\text{muscle}} + J_c^T(\bm{q}) \bm{F}_{\text{contact}},
  \label{eq:ch5:eom-bio}
\end{equation}
where:
\begin{itemize}
  \item $M(\bm{q}) \in \Reals^{n \times n}$ is the mass matrix (from segment inertial properties)
  \item $C(\bm{q}, \dot{\bm{q}}) \in \Reals^{n \times n}$ is the Coriolis/centrifugal matrix
  \item $g(\bm{q}) \in \Reals^n$ is the gravity vector
  \item $R(\bm{q}) \in \Reals^{n \times p}$ is the muscle moment arm matrix ($p$ muscles, $n$ DOF)
  \item $\bm{F}_{\text{muscle}} \in \Reals^p$ are the muscle forces (from Hill models, Chapter~3)
  \item $J_c(\bm{q})$ is the contact Jacobian
  \item $\bm{F}_{\text{contact}}$ are external contact forces
\end{itemize}

The key difference from Volume~I's $\bm{\tau} = B\bm{u}$ is that joint torques are
no longer direct control inputs. Instead:
\begin{equation}
  \bm{\tau}_{\text{net}} = R(\bm{q}) \bm{F}_{\text{muscle}}
  = \sum_{i=1}^{p} \bm{r}_i(\bm{q}) \cdot F_{\text{muscle},i},
  \label{eq:ch5:net-torque}
\end{equation}
where $\bm{r}_i(\bm{q}) \in \Reals^n$ is the $i$-th column of $R$, giving the moment
arm vector of muscle $i$ across all joints.

\subsection{Bi-articular Muscles}

Unlike engineering actuators that span a single joint, many muscles cross two
or more joints. These \textbf{bi-articular} (or multi-articular) muscles create
mechanical coupling between joints that has no analog in standard robotics.

\begin{principle}{The Coordination Function of Bi-articular Muscles}{biarticular}
  Bi-articular muscles serve two distinct functions:
  \begin{enumerate}
    \item \textbf{Energy transfer}: they transport mechanical energy from one joint to
          another without metabolic cost. When the proximal joint extends (muscle
          shortening at proximal end) and the distal joint flexes (muscle lengthening
          at distal end), energy is transferred proximally to distally.
    \item \textbf{Coordination constraint}: they link joint motions, reducing the
          effective degrees of freedom and simplifying the control problem.
  \end{enumerate}
  This is a form of \emph{mechanical intelligence}---the morphology of the
  musculoskeletal system performs part of the control task. The gastrocnemius,
  for example, crosses both the knee and ankle, coupling knee extension to ankle
  plantarflexion during the push-off phase of gait.

  \noindent
  \textbf{Control-Theoretic Perspective:} In the language of Agrachev \& Sachkov's geometric
  control theory \cite{AgrachevSachkov2004}, the moment arm matrix $R(\bm{q})$ is a covector
  field describing how muscle forces map to joint torques. Bi-articular muscles induce dependencies
  in the columns of $R(\bm{q})$, effectively constraining the input distribution $G(\bm{x})$ of the
  system. This reduces the effective control dimension and is formalized in Chapter~3 of
  \cite{AgrachevSachkov2004} (reachability and controllability of affine systems).
\end{principle}

\begin{lstlisting}[language=Python, caption={Computing joint torques from muscle forces through moment arms.}]
import numpy as np

def compute_joint_torques(
    moment_arms: np.ndarray,
    muscle_forces: np.ndarray
) -> np.ndarray:
    """Compute net joint torques from muscle forces and moment arms.

    Parameters
    ----------
    moment_arms : np.ndarray, shape (n_joints, n_muscles)
        Moment arm matrix R(q). Each column gives the moment arm
        of one muscle across all joints. Sign convention: positive
        moment arm produces positive joint torque (typically flexion).
    muscle_forces : np.ndarray, shape (n_muscles,)
        Force produced by each muscle (always positive).

    Returns
    -------
    np.ndarray, shape (n_joints,)
        Net joint torques.
    """
    return moment_arms @ muscle_forces


# Example: 2-DOF elbow system with 4 muscles
# Joints: shoulder flexion, elbow flexion
# Muscles: anterior deltoid, biceps (biarticular),
#          triceps (biarticular), brachialis
#
# Convention: R has shape (n_joints, n_muscles) so that
# tau = R @ F maps the muscle-force vector to a joint-torque
# vector via a single matrix-vector product. Each row is a joint
# and each column is a muscle. This is the shape that
# compute_joint_torques() expects.
R = np.array([
    # deltoid, biceps, triceps, brachialis
    [0.05,  0.03,  -0.025,  0.0 ],  # row 0: shoulder moment arms (m)
    [0.0,   0.04,  -0.020,  0.03],  # row 1: elbow moment arms (m)
])  # Shape: (2 joints, 4 muscles) -- ready to use as-is

# Note: biceps has moment arms at BOTH joints (bi-articular)
muscle_forces = np.array([200.0, 150.0, 100.0, 120.0])  # Newtons

tau = compute_joint_torques(R, muscle_forces)
print(f"Shoulder torque: {tau[0]:.1f} N*m")
print(f"Elbow torque:    {tau[1]:.1f} N*m")
\end{lstlisting}

\subsection{Co-Contraction and Joint Stiffness}

When antagonist muscles co-contract (both agonist and antagonist are active
simultaneously), the net joint torque may be zero but the joint \textbf{stiffness}
increases. This is because each active muscle acts as a spring:
\begin{equation}
  K_{\text{joint}} = \sum_{i=1}^{p} r_i^2(\bm{q}) \cdot \frac{\partial F_{\text{muscle},i}}{\partial \ell_i^M},
  \label{eq:ch5:joint-stiffness}
\end{equation}
where $\partial F / \partial \ell^M$ is the short-range stiffness of the muscle, which
depends on activation level.

Co-contraction is metabolically expensive but biomechanically important: it allows the
nervous system to modulate joint \emph{impedance} independently of joint
\emph{torque}. This is a form of \textbf{variable impedance control} that has no direct
analog in rigid robots with torque-controlled joints.

Position-dependent stiffness modulation is a nonlinear phenomenon central to the geometric
analysis of affine control systems. Agrachev \& Sachkov \cite{AgrachevSachkov2004} develop
the theory of feedback linearization and stabilization via nonlinear drift terms (Chapter~6),
which encompasses exactly this scenario: using position-dependent actuation to achieve
desired closed-loop dynamics.

\section{The Mass Matrix of the Human Body}

Computing the mass matrix $M(\bm{q})$ for a full musculoskeletal model follows the
same recursive algorithms as Volume~0, Chapter~10 (the Composite Rigid Body Algorithm),
using the segment inertial properties from Chapter~2. The key computational steps are:

\begin{enumerate}
  \item Define the kinematic tree (segment connectivity and joint types)
  \item Compute forward kinematics using the product of exponentials
  \item Apply the CRBA to compute $M(\bm{q})$
\end{enumerate}

For a 20-DOF model (pelvis + 2 legs + trunk + 2 arms), $M(\bm{q})$ is a
$20 \times 20$ symmetric positive-definite matrix.

\section*{Exercises}

\begin{enumerate}
  \item \textbf{Bi-articular Energy Transfer.}
        For the gastrocnemius (crossing knee and ankle), show mathematically how
        knee extension can produce ankle torque through the bi-articular coupling.

  \item \textbf{Co-Contraction Analysis.}
        Compute the joint stiffness when biceps and triceps co-contract at 50\%
        activation with the elbow at $90^\circ$.

  \item \textbf{Mass Matrix Properties.}
        For a 2-segment arm model, compute $M(\bm{q})$ at three configurations and
        verify positive-definiteness. Plot the condition number as a function of elbow angle.

  \item \textbf{(Programming.)} Implement the CRBA for a 3-segment leg model using
        the de~Leva segment properties. Compare with Pinocchio output.

  \item \textbf{(Programming.)} Simulate a 2-DOF arm with 4 muscles (as in the code
        listing) performing a reaching movement. Use prescribed muscle activations
        and compute the resulting trajectory via forward dynamics.
\end{enumerate}

\chapter{Inverse Problems in Biomechanics}
\label{ch:inverse-problems}

\epigraph{%
  In biomechanics, we almost never measure forces directly.
  We measure motion and infer the forces that caused it.
}{}

\section{The Inverse Dynamics Pipeline}

Inverse dynamics is the most widely used computational tool in biomechanics: given
measured kinematics $\bm{q}(t)$, $\dot{\bm{q}}(t)$, $\ddot{\bm{q}}(t)$ and external
forces $\bm{F}_{\text{ext}}(t)$, compute the net joint torques. The complete pipeline,
illustrated in Figure~\ref{fig:ch6:id-pipeline}, transforms raw experimental data into
mechanistic estimates of the forces acting on the system:
\begin{equation}
  \bm{\tau}(t) = M(\bm{q})\ddot{\bm{q}} + C(\bm{q}, \dot{\bm{q}})\dot{\bm{q}} + g(\bm{q})
  - J_c^T(\bm{q}) \bm{F}_{\text{ext}}.
  \label{eq:ch6:inverse-dynamics}
\end{equation}

\begin{center}
\begin{tikzpicture}[scale=0.9, >=Stealth, node distance=2.2cm]
  \node[draw, thick, minimum width=2.5cm, fill=blue!10] (collect) {Data Collection\\(markers, forces)};
  \node[draw, thick, minimum width=2.5cm, fill=blue!10, right of=collect] (process) {Marker\\Processing};
  \node[draw, thick, minimum width=2.5cm, fill=green!10, right of=process] (ik) {Inverse\\Kinematics};
  \node[draw, thick, minimum width=2.5cm, fill=green!10, right of=ik] (id) {Inverse\\Dynamics};
  \node[draw, thick, minimum width=2.5cm, fill=red!10, right of=id] (muscle) {Muscle Force\\Estimation};

  \draw[->, thick] (collect) -- (process);
  \draw[->, thick] (process) -- (ik);
  \draw[->, thick] (ik) -- (id);
  \draw[->, thick] (id) -- (muscle);

  \node[below, font=\small] at (collect.south) {Observables};
  \node[below, font=\small] at (ik.south) {Kinematics};
  \node[below, font=\small] at (id.south) {Net torques};
  \node[below, font=\small] at (muscle.south) {Individual forces};
\end{tikzpicture}
\end{center}

The standard pipeline comprises five stages:
\begin{enumerate}
  \item \textbf{Data Collection}: Motion capture markers + force plates
  \item \textbf{Marker Processing}: Filtering, gap-filling, labeling
  \item \textbf{Inverse Kinematics}: Fit skeleton model to marker data
  \item \textbf{Inverse Dynamics}: Compute joint torques from RNEA
  \item \textbf{Muscle Force Estimation}: Distribute torques among muscles
\end{enumerate}

\subsection{Inverse Kinematics}

Given marker positions $\bm{m}_i^{\text{meas}}(t) \in \Reals^3$ for $i = 1, \ldots, N_m$
markers, inverse kinematics solves:
\begin{equation}
  \bm{q}^*(t) = \argmin_{\bm{q}} \sum_{i=1}^{N_m} w_i
  \norm{\bm{m}_i^{\text{model}}(\bm{q}) - \bm{m}_i^{\text{meas}}(t)}^2,
  \label{eq:ch6:ik-optimization}
\end{equation}
where $\bm{m}_i^{\text{model}}(\bm{q})$ is the predicted marker position from the
skeleton model and $w_i$ are weighting factors.

\begin{lstlisting}[language=Python, caption={Simplified inverse kinematics solver.}]
import numpy as np
from scipy.optimize import minimize

def inverse_kinematics_frame(
    marker_positions_measured: np.ndarray,
    marker_positions_model_func,
    q_initial: np.ndarray,
    weights: np.ndarray | None = None
) -> np.ndarray:
    """Solve inverse kinematics for a single frame.

    Parameters
    ----------
    marker_positions_measured : np.ndarray, shape (n_markers, 3)
        Measured marker positions.
    marker_positions_model_func : callable
        Function q -> (n_markers, 3) predicted marker positions.
    q_initial : np.ndarray, shape (n_dof,)
        Initial guess for joint angles.
    weights : np.ndarray, shape (n_markers,), optional
        Per-marker weights.

    Returns
    -------
    np.ndarray, shape (n_dof,)
        Optimal joint angles.
    """
    n_markers = marker_positions_measured.shape[0]
    if weights is None:
        weights = np.ones(n_markers)

    def objective(q: np.ndarray) -> float:
        m_model = marker_positions_model_func(q)
        residuals = m_model - marker_positions_measured
        return float(np.sum(weights[:, None] * residuals**2))

    result = minimize(objective, q_initial, method='L-BFGS-B')
    return result.x
\end{lstlisting}

\subsection{Bottom-Up vs.\ Top-Down Inverse Dynamics}

\textbf{Bottom-up}: start from the ground contact (where forces are measured by
force plates) and propagate upward through the kinematic chain. Most accurate when
force plate data is available.

\textbf{Top-down}: start from the most distal segment and work inward. Useful when
ground reaction forces are not available (e.g., swimming, throwing).

\section{The Muscle Redundancy Problem}

Inverse dynamics gives net joint torques $\bm{\tau} \in \Reals^n$, but we want
individual muscle forces $\bm{F} \in \Reals^p$ with $p > n$ (typically $p = 3n$
to $5n$). This is an underdetermined system:
\begin{equation}
  R(\bm{q}) \bm{F} = \bm{\tau}, \quad \bm{F} \geq 0.
  \label{eq:ch6:redundancy}
\end{equation}

\subsection{Static Optimization}

The most common approach minimizes a cost function:
\begin{equation}
  \min_{\bm{F}} \sum_{i=1}^{p} \left(\frac{F_i}{F_{0,i}}\right)^\alpha
  \quad \text{subject to} \quad R(\bm{q})\bm{F} = \bm{\tau},\;
  0 \leq F_i \leq a_i \cdot f_L(\tilde{\ell}_i) \cdot F_{0,i},
  \label{eq:ch6:static-opt}
\end{equation}
where $\alpha = 2$ (minimize sum of squared activations) or $\alpha = 3$ (minimize
metabolic energy proxy). The constraint $F_i \leq a_i \cdot f_L \cdot F_{0,i}$ enforces
the force-length and activation constraints from Chapter~3.

\subsection{Computed Muscle Control (CMC)}

CMC \cite{Thelen2003} combines forward dynamics with feedback control. It uses a
PD controller to track the measured kinematics, with the controller output being
the desired muscle excitations:
\begin{equation}
  \ddot{\bm{q}}_{\text{des}} = \ddot{\bm{q}}_{\text{meas}} + K_v(\dot{\bm{q}}_{\text{meas}} - \dot{\bm{q}})
  + K_p(\bm{q}_{\text{meas}} - \bm{q}),
  \label{eq:ch6:cmc}
\end{equation}
where $K_v$ and $K_p$ are velocity and position feedback gains.

\section{Parameter Estimation}

Many biomechanical parameters (segment masses, muscle strengths, joint axis locations)
are not directly measurable and must be estimated from motion data.

\subsection{Bayesian Framework}

The parameter estimation problem is naturally Bayesian:
\begin{equation}
  p(\bm{\theta} | \bm{y}) \propto p(\bm{y} | \bm{\theta}) \cdot p(\bm{\theta}),
  \label{eq:ch6:bayes}
\end{equation}
where $\bm{\theta}$ are model parameters, $\bm{y}$ are observations, $p(\bm{y} | \bm{\theta})$
is the likelihood, and $p(\bm{\theta})$ is the prior (from literature values and
physical constraints).

\section*{Exercises}

\begin{enumerate}
  \item \textbf{Inverse Dynamics.}
        For a 2-segment pendulum model of the arm, compute the shoulder and elbow
        torques required during a reaching movement given $\bm{q}(t) = (\sin(\pi t/T), \pi t / (2T))$.

  \item \textbf{Muscle Redundancy.}
        For a single-DOF elbow with 3 muscles (biceps, brachialis, brachioradialis),
        solve the static optimization for a given elbow torque. Show that the solution
        depends on the cost function exponent~$\alpha$.

  \item \textbf{(Programming.)} Implement the static optimization for a 2-DOF,
        6-muscle system. Compare $\alpha = 2$ and $\alpha = 3$ solutions.

  \item \textbf{(Programming.)} Implement an inverse kinematics solver for a
        3-segment arm model with 7 markers. Test with synthetic marker data.
\end{enumerate}

\chapter{Experimental Methods}
\label{ch:experimental}

\epigraph{%
  The quality of a biomechanical analysis is limited by the quality of the measurements,
  not the sophistication of the model.
}{}

\section{Motion Capture Systems}

\subsection{Marker-Based Systems}

Optoelectronic motion capture systems (Vicon, Qualisys, OptiTrack) track retroreflective
markers attached to the body surface. Standard configurations use 30--80 markers placed
on bony landmarks following established protocols (Plug-in-Gait, Cleveland Clinic, ISB).

Key specifications for biomechanical applications
\cite{chiari2005human,chiari2005human}:
\begin{itemize}
  \item \textbf{Sampling rate}: 100--500~Hz for human motion (higher for impacts)
  \item \textbf{Spatial accuracy}: $<1$~mm in calibrated volume
  \item \textbf{Marker size}: 9--14~mm diameter
  \item \textbf{Latency}: important for real-time applications ($<10$~ms achievable)
\end{itemize}

It is important to note that the instrument accuracy above refers to the camera
system alone. The dominant error in marker-based biomechanics is not camera noise
but \emph{soft tissue artifact}---the relative motion between skin markers and the
underlying bone---discussed in Section~\ref{sec:sta} below
\cite{chiari2005human,chiari2005human}.

\subsection{Markerless Systems}

Computer vision approaches (OpenPose, MediaPipe, DeepLabCut) estimate joint positions
from video without physical markers. Accuracy is currently 10--30~mm, insufficient for
clinical-grade inverse dynamics but useful for field measurements and large-scale studies.

\begin{lstlisting}[language=Python, caption={Motion capture data processing pipeline.}]
import numpy as np
from scipy.signal import butter, filtfilt

def process_mocap_data(
    marker_data: np.ndarray,
    sample_rate: float,
    cutoff_freq: float = 6.0,
    filter_order: int = 4
) -> tuple[np.ndarray, np.ndarray, np.ndarray]:
    """Process raw motion capture marker data.

    Parameters
    ----------
    marker_data : np.ndarray, shape (n_frames, n_markers, 3)
        Raw marker positions in meters.
    sample_rate : float
        Sampling frequency in Hz.
    cutoff_freq : float
        Low-pass filter cutoff frequency in Hz.
    filter_order : int
        Butterworth filter order.

    Returns
    -------
    tuple of np.ndarray
        (positions, velocities, accelerations)
        Each shape (n_frames, n_markers, 3).
    """
    # Design low-pass Butterworth filter
    nyquist = sample_rate / 2.0
    b, a = butter(filter_order, cutoff_freq / nyquist, btype='low')

    n_frames, n_markers, _ = marker_data.shape
    positions = np.zeros_like(marker_data)
    velocities = np.zeros_like(marker_data)
    accelerations = np.zeros_like(marker_data)

    dt = 1.0 / sample_rate

    for m in range(n_markers):
        for axis in range(3):
            # Filter position data
            positions[:, m, axis] = filtfilt(
                b, a, marker_data[:, m, axis]
            )

    # Central difference for velocities
    velocities[1:-1] = (positions[2:] - positions[:-2]) / (2 * dt)
    velocities[0] = velocities[1]
    velocities[-1] = velocities[-2]

    # Central difference for accelerations
    accelerations[1:-1] = (
        positions[2:] - 2 * positions[1:-1] + positions[:-2]
    ) / dt**2
    accelerations[0] = accelerations[1]
    accelerations[-1] = accelerations[-2]

    return positions, velocities, accelerations
\end{lstlisting}

\subsection{Soft Tissue Artifact}
\label{sec:sta}

Soft tissue artifact (STA) is the relative motion between skin-mounted markers
and the underlying bone caused by deformation of skin, subcutaneous fat, and muscle
during movement. STA is widely recognized as the dominant source of error in
marker-based motion capture, frequently exceeding the nominal camera resolution
by more than an order of magnitude \cite{chiari2005human,peters2010quantification}.

The magnitude of STA varies strongly with marker location, anatomical region,
task, and subject anthropometry. Representative values for the lower limb are
10--30~mm of skin-marker displacement relative to bone during walking and running,
with the thigh segment the worst offender and the shank the best
\cite{chiari2005human,peters2010quantification}. STA is not random noise: it is
task-correlated and heavily dependent on muscle contraction state, which means
low-pass filtering cannot remove it.

Practical mitigation strategies include:
\begin{itemize}
  \item \textbf{Rigid marker clusters} mounted on lightweight shells, which average
        out local skin motion over a neighborhood of markers.
  \item \textbf{Optimal pose estimation}, in which the rigid body pose of a segment
        is obtained by least-squares fitting of a cluster of markers rather than by
        differencing pairs of landmarks \cite{chiari2005human}.
  \item \textbf{Global optimization / inverse kinematics} with joint constraints,
        which enforces an articulated model and redistributes STA across the chain.
  \item \textbf{Biplanar fluoroscopy} or bone pins when the research question
        demands sub-millimeter bone-to-bone accuracy.
\end{itemize}

STA should always be reported as part of the uncertainty budget of any kinetic
analysis; inverse dynamics is sensitive to segment orientation errors, and
moderate STA can produce large errors in computed joint moments.

\subsection{Low-Pass Filter Cutoff Selection}
\label{sec:filter-cutoff}

Choosing the cutoff frequency for the low-pass Butterworth filter in
\texttt{process\_mocap\_data()} is not a matter of taste: a cutoff that is too low
attenuates the true signal (especially accelerations, whose energy extends to
higher frequencies than positions), while a cutoff that is too high leaves noise
that is amplified by numerical differentiation.

The standard objective procedure is \textbf{residual analysis}, due to Winter
\cite{Winter2009}. For a range of candidate cutoffs $f_c$, one computes
the root-mean-square residual
\[
  R(f_c) = \sqrt{\frac{1}{N}\sum_{i=1}^{N}\bigl(x_i - \hat{x}_i(f_c)\bigr)^2},
\]
where $x_i$ is the raw signal and $\hat{x}_i(f_c)$ is the filtered signal. Plotting
$R(f_c)$ vs.\ $f_c$ produces a curve that is approximately linear in the high-$f_c$
regime (residuals reflect only noise) and rises sharply at low $f_c$ (residuals
include signal). The optimal cutoff is the intersection of the extrapolated linear
noise-only segment with the $R$-axis, typically in the range 4--10~Hz for
kinematics of whole-body locomotion and 10--20~Hz for running and impact-dominated
tasks \cite{Winter2009}. Higher cutoffs are required when downstream
analysis needs accelerations or joint powers, since differentiation acts as a
high-pass amplifier.

\section{Electromyography (EMG)}

EMG measures the electrical activity of muscles, providing a window into neural
commands. Surface EMG (sEMG) is non-invasive; intramuscular EMG (using fine-wire
or needle electrodes) targets deep muscles.

\subsection{EMG Processing}

The raw EMG signal must be processed before use:
\begin{enumerate}
  \item \textbf{Band-pass filter}: 20--450~Hz (remove movement artifact and noise)
  \item \textbf{Full-wave rectification}: $|e(t)|$
  \item \textbf{Envelope extraction}: low-pass filter at 3--6~Hz
  \item \textbf{Normalization}: divide by maximum voluntary contraction (MVC) value
\end{enumerate}

\section{Force Plates}

Force plates measure ground reaction forces (GRF) and moments. Standard specifications:
\begin{itemize}
  \item 6-component measurement: $F_x, F_y, F_z, M_x, M_y, M_z$
  \item Sampling rate: 1000--2000~Hz
  \item Range: 0--5000~N vertical, 0--2500~N horizontal
  \item Center of pressure (CoP) accuracy: $<2$~mm
\end{itemize}

The center of pressure is computed from force plate outputs:
\begin{equation}
  x_{\text{CoP}} = -\frac{M_y + F_x \cdot d_z}{F_z}, \quad
  y_{\text{CoP}} = \frac{M_x - F_y \cdot d_z}{F_z},
  \label{eq:ch7:cop}
\end{equation}
where $d_z$ is the distance from the force plate surface to the transducer plane.

\section{Inertial Measurement Units (IMUs)}

IMUs combine accelerometers, gyroscopes, and magnetometers in a compact package.
They provide:
\begin{itemize}
  \item Angular velocity: $\bm{\omega}(t) \in \Reals^3$
  \item Linear acceleration: $\bm{a}(t) \in \Reals^3$
  \item Magnetic field: $\bm{B}(t) \in \Reals^3$ (for heading)
\end{itemize}

Orientation estimation from IMU data is a state estimation problem solvable using
the extended Kalman filter (Volume~I, Chapter~6) or complementary filters.

\section{Data Synchronization}

Multiple measurement systems must be temporally synchronized. Best practices:
\begin{itemize}
  \item Hardware sync: trigger all systems from a common clock pulse
  \item Minimum: force plate and motion capture on the same clock
  \item EMG-to-motion delay: typically 10--80~ms (electromechanical delay)
  \item Cross-correlation for post-hoc alignment when hardware sync is unavailable
\end{itemize}

\section*{Exercises}

\begin{enumerate}
  \item \textbf{Filter Selection.}
        Using the provided code, filter a synthetic marker trajectory with cutoff
        frequencies of 4, 6, 8, and 12~Hz. Plot the effect on position, velocity,
        and acceleration. Discuss the trade-off between noise reduction and signal distortion.

  \item \textbf{CoP Computation.}
        Given force plate data from a vertical jump ($F_z$ peak = 3000~N), compute
        the center of pressure trajectory during takeoff.

  \item \textbf{EMG Processing.}
        Process a synthetic EMG signal (sum of sinusoids + white noise) through the
        full processing pipeline. Plot each stage.

  \item \textbf{(Programming.)} Implement an IMU orientation estimator using a
        complementary filter. Test with synthetic gyroscope and accelerometer data.
\end{enumerate}

\chapter{Inference on Biological Systems}
\label{ch:inference}

\begin{gomdraftnotice}
This chapter is under active development and contains only a structural outline
with preliminary material. Full exposition, worked examples, proofs, and exercises
will be added in future editions.
\end{gomdraftnotice}

\epigraph{%
  We cannot open the body and read the parameters.
  We must infer them from the traces left by motion.
}{}

\section{Parameter Identification from Motion Data}

Biological system parameters (muscle strengths, tendon slack lengths, segment inertias)
cannot typically be measured directly in vivo. They must be inferred from observable
motion data using optimization and estimation techniques.

\begin{center}
\begin{tikzpicture}[scale=0.9, >=Stealth, node distance=2cm]
  \node[draw, thick, fill=blue!10, minimum width=2cm] (model) {Model $f(\state, \bm{\theta})$};
  \node[draw, thick, fill=blue!10, minimum width=2cm, right of=model] (predict) {Prediction $\hat{\bm{y}}$};
  \node[draw, thick, fill=green!10, minimum width=2cm, align=center, right of=predict] (compare) {Compare\\to data};
  \node[draw, thick, fill=red!10, minimum width=2cm, below of=compare] (optimize) {Optimize $\bm{\theta}$};

  \node[left, font=\tiny] at (model.west) {Parameters};
  \node[below, font=\tiny] at (predict.south) {Observations};

  \draw[->, thick] (model) -- (predict);
  \draw[->, thick] (predict) -- (compare);
  \draw[->, thick] (compare) -- (optimize);
  \draw[<->, thick] (optimize) -- (model);

  \node[draw, thick, fill=orange!10, minimum width=2cm, below of=model] (prior) {Prior $p(\bm{\theta})$};
  \draw[->, dashed] (prior) -- (optimize) node[midway, left, font=\tiny] {Bayes};
\end{tikzpicture}
\end{center}

\subsection{Maximum Likelihood Estimation}

Given a parameterized model $\dot{\state} = f(\state, \bm{\theta})$ with parameters
$\bm{\theta}$ and noisy observations $\bm{y}_k = h(\state_k) + \bm{v}_k$ at times
$t_k$, the maximum likelihood estimate is:
\begin{equation}
  \hat{\bm{\theta}} = \argmin_{\bm{\theta}} \sum_{k=1}^{N}
  \norm{\bm{y}_k - h(\state_k(\bm{\theta}))}_{R^{-1}}^2,
  \label{eq:ch8:mle}
\end{equation}
where $R$ is the measurement noise covariance.

\subsection{Bayesian Estimation with Physical Constraints}

The Bayesian framework (introduced in Chapter~6) is particularly valuable for
biomechanical parameter estimation because it naturally incorporates prior knowledge
from the literature:

\begin{lstlisting}[language=Python, caption={Bayesian parameter estimation for musculoskeletal model.}]
import numpy as np
from scipy.optimize import minimize

def log_posterior(
    theta: np.ndarray,
    observations: np.ndarray,
    forward_model,
    prior_mean: np.ndarray,
    prior_cov: np.ndarray,
    noise_cov: np.ndarray
) -> float:
    """Compute the log-posterior for parameter estimation.

    Parameters
    ----------
    theta : np.ndarray
        Model parameters.
    observations : np.ndarray, shape (n_obs, n_dims)
        Observed data.
    forward_model : callable
        theta -> predicted_observations.
    prior_mean : np.ndarray
        Prior parameter mean.
    prior_cov : np.ndarray
        Prior parameter covariance.
    noise_cov : np.ndarray
        Observation noise covariance.

    Returns
    -------
    float
        Negative log-posterior (for minimization).
    """
    # Log-likelihood
    predicted = forward_model(theta)
    residuals = observations - predicted
    R_inv = np.linalg.inv(noise_cov)
    log_likelihood = -0.5 * np.sum(residuals @ R_inv * residuals)

    # Log-prior
    prior_residuals = theta - prior_mean
    P_inv = np.linalg.inv(prior_cov)
    log_prior = -0.5 * prior_residuals @ P_inv @ prior_residuals

    return -(log_likelihood + log_prior)


def map_estimate(
    observations: np.ndarray,
    forward_model,
    prior_mean: np.ndarray,
    prior_cov: np.ndarray,
    noise_cov: np.ndarray,
    bounds: list[tuple[float, float]] | None = None
) -> np.ndarray:
    """Compute Maximum A Posteriori (MAP) parameter estimate.

    Parameters
    ----------
    observations : np.ndarray
        Observed data.
    forward_model : callable
        Parameter-to-observation mapping.
    prior_mean : np.ndarray
        Prior mean from literature.
    prior_cov : np.ndarray
        Prior uncertainty.
    noise_cov : np.ndarray
        Measurement noise covariance.
    bounds : list of tuple, optional
        Physical bounds on parameters.

    Returns
    -------
    np.ndarray
        MAP parameter estimate.
    """
    result = minimize(
        log_posterior,
        prior_mean,
        args=(observations, forward_model, prior_mean,
              prior_cov, noise_cov),
        method='L-BFGS-B',
        bounds=bounds
    )
    return result.x
\end{lstlisting}

\section{System Identification for Biological Systems}

System identification techniques from control theory (Volume~I) apply to biomechanical
systems with modifications:

\begin{enumerate}
  \item \textbf{Joint impedance identification}: Perturb the limb and measure the
        response. The stiffness, damping, and inertia of the joint can be estimated
        from the transfer function.
  \item \textbf{Muscle parameter identification}: Optimize muscle parameters (F0,
        optimal length, tendon slack length) to match experimental force or torque data.
  \item \textbf{Neural control identification}: Estimate the feedback gains used by
        the nervous system from perturbation experiments.
\end{enumerate}

\section{Machine Learning Approaches}

Modern ML approaches to biomechanical inference include:
\begin{itemize}
  \item Physics-informed neural networks (PINNs) that enforce equations of motion
  \item Gaussian process regression for joint angle estimation from sparse data
  \item Deep learning for markerless motion capture pose estimation
\end{itemize}

\section*{Exercises}

\begin{enumerate}
  \item \textbf{MAP Estimation.}
        Using the provided code, estimate the optimal fiber length of a biceps muscle
        from simulated isometric force data at 5 joint angles.

  \item \textbf{Identifiability Analysis.}
        For a Hill-type muscle model, analyze which parameters can be uniquely estimated
        from isometric force-angle data alone. Which require dynamic data?

  \item \textbf{(Programming.)} Implement a joint impedance identification experiment:
        simulate a perturbed pendulum and estimate stiffness and damping from the response.
\end{enumerate}

\chapter{Deformable Bodies and Continuum Mechanics}
\label{ch:deformable}

\begin{gomdraftnotice}
This chapter is under active development and contains only a structural outline
with preliminary material. Full exposition, worked examples, proofs, and exercises
will be added in future editions.
\end{gomdraftnotice}

\epigraph{%
  When the rigid-body assumption fails, we do not abandon mechanics.
  We embrace the continuum.
}{}

\section{When Rigid-Body Assumptions Fail}

The rigid-body model fails when:
\begin{enumerate}
  \item \textbf{Internal stresses are needed}: to predict injury, fracture, or tissue damage
  \item \textbf{Deformation affects kinematics}: shaft flexibility in golf, spine compliance
  \item \textbf{Wave propagation matters}: impact biomechanics, blast injury
  \item \textbf{Fluid-structure interaction}: blood flow, synovial fluid, respiratory mechanics
\end{enumerate}

\section{Beam Models for Long Bones and Shafts}

When deformation is small and the body is slender (length $\gg$ diameter), beam
theory provides an efficient model. A schematic representation of the finite element
beam model is shown in Figure~\ref{fig:ch9:beam-fem}.

\begin{figure}[htbp]
\centering
\begin{tikzpicture}[scale=0.85, >=Stealth]
  % Continuous beam
  \draw[very thick] (0, 2) -- (6, 2.3) node[right] {Euler-Bernoulli beam};
  \filldraw (0, 2) circle (3pt);
  \node[below] at (0, 2) {Clamped};

  % Distributed load
  \draw[->, thick, red] (1.5, 2.8) -- (1.5, 2.2) node[right] {$f(x,t)$};
  \draw[->, thick, red] (3, 2.9) -- (3, 2.2);
  \draw[->, thick, red] (4.5, 2.8) -- (4.5, 2.2);

  % Finite elements
  \node at (3, 0.8) {\textbf{Discretization into FE elements}};
  \filldraw (0, 0) circle (2pt);
  \filldraw (1.5, 0.1) circle (2pt);
  \filldraw (3, 0) circle (2pt);
  \filldraw (4.5, 0.05) circle (2pt);
  \filldraw (6, 0.2) circle (2pt);

  \draw (0, 0) -- (1.5, 0.1) -- (3, 0) -- (4.5, 0.05) -- (6, 0.2);

  % Deformation
  \draw[dashed, gray] (0, 2) -- (0, 0);
  \draw[dashed, gray] (6, 2.3) -- (6, 0.2);
\end{tikzpicture}
\caption{Finite element beam model: continuous Euler-Bernoulli beam (top) and its discretization into nodal elements (bottom).}
\label{fig:ch9:beam-fem}
\end{figure}

For a Euler-Bernoulli beam:
\begin{equation}
  EI \frac{\partial^4 w}{\partial x^4} + \rho A \frac{\partial^2 w}{\partial t^2} = f(x, t),
  \label{eq:ch9:euler-bernoulli}
\end{equation}
where $E$ is the elastic modulus, $I$ is the second moment of area, $\rho$ is
density, $A$ is cross-sectional area, $w(x,t)$ is transverse displacement, and
$f$ is the distributed load.

\subsection{Modal Analysis}

The solution is expanded in modes:
\begin{equation}
  w(x, t) = \sum_{k=1}^{N} \phi_k(x) \eta_k(t),
  \label{eq:ch9:modal-expansion}
\end{equation}
where $\phi_k(x)$ are mode shapes and $\eta_k(t)$ are modal coordinates. This converts
the PDE to a system of ODEs that can be appended to the rigid-body equations of motion.

\subsection{Golf Club Shaft Flexibility}

The golf club shaft is the most accessible example of a flexible link in sport
biomechanics. A typical shaft has:
\begin{itemize}
  \item Length: $\sim$1.1~m
  \item First bending mode: 4--6~Hz (driver), 6--10~Hz (irons)
  \item Tip deflection at impact: 5--15~cm
  \item Lead/lag deflection: 2--5~cm
  \item Toe-down droop: 3--8~cm (gravity effect during downswing)
\end{itemize}

\begin{lstlisting}[language=Python, caption={Modal analysis of a golf club shaft.}]
import numpy as np
from scipy.linalg import eigh

def shaft_modal_analysis(
    n_elements: int = 20,
    length: float = 1.1,
    EI_root: float = 50.0,
    EI_tip: float = 15.0,
    rho_A: float = 0.15
) -> tuple[np.ndarray, np.ndarray]:
    """Compute natural frequencies and mode shapes of a tapered shaft.

    Uses finite element beam model with linear EI taper.

    Parameters
    ----------
    n_elements : int
        Number of finite elements.
    length : float
        Shaft length in meters.
    EI_root : float
        Flexural rigidity at root (grip end) in N*m^2.
    EI_tip : float
        Flexural rigidity at tip (head end) in N*m^2.
    rho_A : float
        Mass per unit length in kg/m.

    Returns
    -------
    tuple of np.ndarray
        (frequencies_hz, mode_shapes)
    """
    n_nodes = n_elements + 1
    n_dof = 2 * n_nodes  # displacement + rotation at each node
    h = length / n_elements

    K = np.zeros((n_dof, n_dof))
    M = np.zeros((n_dof, n_dof))

    for e in range(n_elements):
        # Local EI for this element (linear taper)
        xi = (e + 0.5) / n_elements
        EI_local = EI_root + (EI_tip - EI_root) * xi

        # Euler-Bernoulli beam element stiffness matrix
        k_e = EI_local / h**3 * np.array([
            [12,   6*h,  -12,   6*h],
            [6*h,  4*h**2, -6*h, 2*h**2],
            [-12, -6*h,   12,  -6*h],
            [6*h,  2*h**2, -6*h, 4*h**2]
        ])

        # Consistent mass matrix
        m_e = rho_A * h / 420 * np.array([
            [156,   22*h,  54,   -13*h],
            [22*h,  4*h**2, 13*h, -3*h**2],
            [54,    13*h,  156,  -22*h],
            [-13*h, -3*h**2, -22*h, 4*h**2]
        ])

        # Assembly
        dofs = [2*e, 2*e+1, 2*e+2, 2*e+3]
        for i_idx, i in enumerate(dofs):
            for j_idx, j in enumerate(dofs):
                K[i, j] += k_e[i_idx, j_idx]
                M[i, j] += m_e[i_idx, j_idx]

    # Boundary conditions: clamped at root (grip)
    free_dofs = list(range(2, n_dof))  # Remove first node
    K_free = K[np.ix_(free_dofs, free_dofs)]
    M_free = M[np.ix_(free_dofs, free_dofs)]

    # Solve eigenvalue problem
    eigenvalues, eigenvectors = eigh(K_free, M_free)

    # Convert to Hz
    frequencies = np.sqrt(np.abs(eigenvalues)) / (2 * np.pi)

    return frequencies[:6], eigenvectors[:, :6]


# Example
freqs, modes = shaft_modal_analysis()
print("Natural frequencies (Hz):")
for i, f in enumerate(freqs):
    print(f"  Mode {i+1}: {f:.1f} Hz")
\end{lstlisting}

\section{Finite Element Methods (Overview)}

For complex geometries (vertebral bodies, skull, articular surfaces), the finite
element method (FEM) provides the most general approach. While a full treatment is
beyond this text's scope, the key concept is:

The continuous domain $\Omega$ is divided into elements, and the displacement field
is approximated using shape functions:
\begin{equation}
  \bm{u}(\bm{x}, t) \approx \sum_{i} N_i(\bm{x}) \bm{u}_i(t),
  \label{eq:ch9:fem-approx}
\end{equation}
leading to the assembled system:
\begin{equation}
  M_{\text{FEM}} \ddot{\bm{u}} + C_{\text{FEM}} \dot{\bm{u}} + K_{\text{FEM}} \bm{u} = \bm{f}_{\text{ext}}.
  \label{eq:ch9:fem-system}
\end{equation}

\section*{Exercises}

\begin{enumerate}
  \item \textbf{Shaft Frequency.}
        Using the provided code, compute how the first natural frequency changes
        as the tip-to-root EI ratio varies from 0.1 to 0.9. Plot the result.

  \item \textbf{Mode Shapes.}
        Plot the first three mode shapes of the shaft. Identify nodes (zero-crossings)
        and antinodes (maximum displacement).

  \item \textbf{Timescale Separation.}
        If the golf downswing takes 0.25~s, estimate the number of oscillation cycles
        for each of the first 3 modes. Which modes are quasi-static?

  \item \textbf{(Programming.)} Add a tip mass (clubhead, 0.2~kg) to the shaft FEM
        model and recompute the natural frequencies. How does the clubhead mass affect
        the frequencies?
\end{enumerate}

\chapter{Applications of Control Theory to Biological Systems}
\label{ch:control-bio}

\epigraph{%
  The body is not a robot that happens to be made of meat.
  It is a system that has evolved to exploit its own physics.
}{}

\section{Forward Dynamics Simulation}

Forward dynamics simulation integrates the equations of motion given muscle
excitations $\bm{u}(t)$:
\begin{equation}
  \ddot{\bm{q}} = M(\bm{q})^{-1}\left[R(\bm{q})\bm{F}_{\text{muscle}} + J_c^T \bm{F}_c
  - C(\bm{q}, \dot{\bm{q}})\dot{\bm{q}} - g(\bm{q})\right],
  \label{eq:ch10:forward-dynamics}
\end{equation}
coupled with the muscle state equations (Chapter~3):
\begin{align}
  \dot{a}_i &= (u_i - a_i) / \tau_i, \\
  \dot{\ell}_i^M &= v_i^M(a_i, \ell_i^M, \ell_i^{MT}(\bm{q})).
\end{align}

\begin{lstlisting}[language=Python, caption={Forward dynamics integration for a musculoskeletal model.}]
import numpy as np
from scipy.integrate import solve_ivp
from typing import Callable

def forward_dynamics_simulation(
    mass_matrix_func: Callable,
    coriolis_func: Callable,
    gravity_func: Callable,
    moment_arm_func: Callable,
    muscle_force_func: Callable,
    excitation_func: Callable,
    q0: np.ndarray,
    qdot0: np.ndarray,
    a0: np.ndarray,
    t_span: tuple[float, float],
    n_eval: int = 500
) -> dict:
    """Simulate forward dynamics of a musculoskeletal system.

    Parameters
    ----------
    mass_matrix_func : callable
        q -> M(q), shape (n_dof, n_dof)
    coriolis_func : callable
        (q, qdot) -> C(q, qdot) @ qdot, shape (n_dof,)
    gravity_func : callable
        q -> g(q), shape (n_dof,)
    moment_arm_func : callable
        q -> R(q), shape (n_dof, n_muscles)
    muscle_force_func : callable
        (a, l_M, v_M) -> F_muscle for each muscle
    excitation_func : callable
        t -> u(t), shape (n_muscles,)
    q0 : np.ndarray
        Initial joint angles.
    qdot0 : np.ndarray
        Initial joint velocities.
    a0 : np.ndarray
        Initial muscle activations.
    t_span : tuple
        (t_start, t_end)
    n_eval : int
        Number of output time points.

    Returns
    -------
    dict
        Keys: 't', 'q', 'qdot', 'a', 'tau'
    """
    n_dof = len(q0)
    n_muscles = len(a0)
    t_eval = np.linspace(t_span[0], t_span[1], n_eval)

    def dynamics(t: float, state: np.ndarray) -> np.ndarray:
        q = state[:n_dof]
        qdot = state[n_dof:2*n_dof]
        a = np.clip(state[2*n_dof:], 0.0, 1.0)

        # Muscle forces (simplified: assume known fiber length)
        F_muscle = np.array([
            a[i] * 500.0  # Simplified: F = a * F0
            for i in range(n_muscles)
        ])

        # Joint torques
        R = moment_arm_func(q)
        tau = R @ F_muscle

        # Skeletal dynamics
        M = mass_matrix_func(q)
        C_qdot = coriolis_func(q, qdot)
        g = gravity_func(q)
        qddot = np.linalg.solve(M, tau - C_qdot - g)

        # Activation dynamics
        u = excitation_func(t)
        tau_act = 0.015  # 15 ms activation
        tau_deact = 0.050  # 50 ms deactivation
        adot = np.where(
            u > a,
            (u - a) / tau_act,
            (u - a) / tau_deact
        )

        return np.concatenate([qdot, qddot, adot])

    state0 = np.concatenate([q0, qdot0, a0])
    sol = solve_ivp(dynamics, t_span, state0, t_eval=t_eval,
                    method='RK45', max_step=0.001)

    return {
        't': sol.t,
        'q': sol.y[:n_dof].T,
        'qdot': sol.y[n_dof:2*n_dof].T,
        'a': sol.y[2*n_dof:].T,
    }
\end{lstlisting}

\section{Predictive Simulation}

Predictive simulation uses trajectory optimization (Volume~I, Chapter~5) to find
the muscle excitations that produce a desired motion while minimizing a cost:
\begin{equation}
  \min_{\bm{u}(\cdot)} \int_0^T \mathcal{L}(\state, \bm{u}) \dt
  \quad \text{s.t.} \quad \dot{\state} = f(\state, \bm{u}),\;
  \state(0) = \state_0,\; h(\state, \bm{u}) \leq 0,
  \label{eq:ch10:predictive-sim}
\end{equation}
where $\mathcal{L}$ is a physiologically motivated cost (e.g., metabolic energy,
muscle fatigue, deviation from desired kinematics) and $h$ captures physical
constraints (joint limits, muscle force capacity, contact constraints).

\noindent
Agrachev \& Sachkov \cite{AgrachevSachkov2004}, Theorem~8.1 (Pontryagin's Maximum Principle),
establishes the fundamental existence and optimality conditions for such optimal control problems.
Their treatment covers affine systems with constraints and provides both theoretical guarantees
and computational methods (the method of characteristics) for solving OCPs of this form.

\subsection{Metabolic Cost Models}

Predictive simulation requires a physiologically meaningful running cost. The
most widely used choice in gait simulation is a \emph{metabolic energy rate} model,
which converts muscle state and force into an estimate of the whole-body
rate of ATP consumption. Umberger et al.\ \cite{umberger2003metabolic} proposed a
four-term decomposition in which the instantaneous metabolic rate of muscle $i$
is written as
\[
  \dot{E}_i = \dot{h}_i^{A} + \dot{h}_i^{M} + \dot{h}_i^{SL} + \dot{w}_i^{CE},
\]
summing activation heat, maintenance heat, shortening/lengthening heat, and the
mechanical work done by the contractile element, each calibrated against in-vitro
muscle energetics. Bhargava et al.\ \cite{bhargava2004metabolic} offer a closely
related formulation that separates activation and maintenance using fast and slow
fiber fractions, and is the model shipped with OpenSim
\cite{delp2007opensim,rajagopal2016full}. Whole-body metabolic cost is obtained by
summing $\dot{E}_i$ across all muscles and adding a constant basal rate.

In practice, minimizing $\int_0^T \sum_i \dot{E}_i \dt$ is expensive and nonsmooth,
so many optimal-control studies substitute a surrogate such as $\sum_i a_i^p$ or
$\sum_i (F_i/F_i^0)^p$ with $p \in \{2, 3\}$
\cite{pandy2001computer,anderson2001dynamic}. Surrogates and calibrated metabolic
models generally agree on qualitative patterns (preferred walking speed, stride
frequency) but can disagree on absolute cost.

\subsection{Muscle Synergies}

Predictive simulation of a 600-muscle model has many more inputs than degrees of
freedom, reproducing the motor-redundancy problem of Chapter~1 at the control
level. A long line of neuroscience research has proposed that the central nervous
system sidesteps this curse of dimensionality by driving muscles in low-dimensional
patterns called \emph{muscle synergies}: a small set of time-invariant activation
vectors $\bm{w}_k \in \Reals^{n_{\text{m}}}$ modulated by time-varying coefficients
$c_k(t)$ so that the full activation is
\[
  \bm{a}(t) \approx \sum_{k=1}^{K} c_k(t)\, \bm{w}_k,
  \qquad K \ll n_{\text{m}}.
\]
Bizzi \& Cheung \cite{bizzi2008combining} summarize the evidence for synergies as
building blocks of motor behavior across species, while d'Avella et al.\
\cite{davella2003combinations} showed that a handful of time-varying synergies
account for most of the EMG variance in natural arm reaching. Non-negative matrix
factorization \cite{HarrisWolpert1998} is the standard extraction technique because
activations and synergies are non-negative.

For control-theoretic purposes, synergies provide a rationale for low-dimensional
parameterizations of the input in optimal-control problems, reducing the effective
control dimension from $n_{\text{m}}$ to $K$ and making predictive simulation
tractable at whole-body scale.

\section{Prosthetics and Exoskeleton Control}

The control theory of Volumes~I--II applies directly to assistive devices:
\begin{itemize}
  \item \textbf{Powered prostheses}: trajectory tracking with impedance modulation
  \item \textbf{Exoskeletons}: assist-as-needed control using contraction theory
  \item \textbf{FES (Functional Electrical Stimulation)}: direct muscle activation
        control for paralyzed limbs
\end{itemize}

\section{Sport Biomechanics Applications}

\subsection{The Golf Swing as a Biomechanical System}

The golf swing demonstrates every concept in this volume:
\begin{enumerate}
  \item Multi-segment kinematic chain (15+ DOF)
  \item Muscle-driven actuation with co-contraction
  \item Bi-articular coupling (forearm muscles, lats)
  \item Passive dynamics exploitation (wrist release)
  \item Shaft flexibility (deformable body)
  \item Ground reaction force coupling (closed chain at the feet)
\end{enumerate}

The complete biomechanical analysis integrates:
\begin{itemize}
  \item Motion capture (kinematics)
  \item Force plates (ground reaction forces)
  \item EMG (muscle activation patterns)
  \item Inverse dynamics (joint torques)
  \item Static optimization (muscle forces)
  \item Forward dynamics validation (predictive simulation)
\end{itemize}

\section{Synthesis: From Theory to Application}

The progression through this volume mirrors the progression through the series:

\begin{center}
\renewcommand{\arraystretch}{1.3}
\begin{tabular}{@{}llll@{}}
\toprule
\textbf{Volume} & \textbf{Concept} & \textbf{Biomechanical Application} & \textbf{Agrachev \& Sachkov} \\
\midrule
Vol~0   & Screw axes, PoE      & Joint kinematics, ISA analysis            & \cite[Ch.~2]{AgrachevSachkov2004} \\
Vol~I   & Contraction theory   & Stability of muscle-driven systems        & \cite[Ch.~4]{AgrachevSachkov2004} \\
Vol~I   & DDP/iLQR             & Optimal muscle excitation patterns        & \cite[Ch.~8]{AgrachevSachkov2004} \\
Vol~II  & Trajectory tubes     & Motor variability envelopes               & \cite[Ch.~3]{AgrachevSachkov2004} \\
Vol~II  & Phase variables      & Gait cycle parameterization               & \cite[Ch.~7]{AgrachevSachkov2004} \\
Vol~III & Hill model           & Force prediction                          & \cite[Ch.~6]{AgrachevSachkov2004} \\
Vol~III & Inverse dynamics     & Joint torque estimation                   & \cite[Ch.~5]{AgrachevSachkov2004} \\
Vol~IV  & Motor control        & Neural command generation                 & \cite[Ch.~8]{AgrachevSachkov2004} \\
\bottomrule
\end{tabular}
\end{center}

\section*{Exercises}

\begin{enumerate}
  \item \textbf{Forward Simulation.}
        Using the provided code (adapted for a simple 2-DOF model), simulate a
        reaching movement with prescribed muscle excitation ramps.

  \item \textbf{Predictive vs.\ Tracking.}
        Compare the muscle activations from static optimization (tracking measured
        kinematics) with those from predictive simulation (optimizing a metabolic cost).

  \item \textbf{Prosthetic Control.}
        Design an impedance controller for a prosthetic knee that provides different
        stiffness during stance and swing phases.

  \item \textbf{(Programming.)} Implement a complete inverse dynamics $\to$ static
        optimization pipeline for a 3-DOF planar arm model with 6 muscles.
\end{enumerate}


\backmatter
\printindex

\bibliographystyle{plain}
\bibliography{../geometry_of_motion}

\end{document}